% !TEX root = ../EDRmain.tex
\def\numb{N}
%\def\manifold{\operatorname{ma}}
\def\manifold{\mathcal{M}}
\def\numbd{\numb^{\circ}}
\def\nmin{\numb_{\operatorname{loc}}}
\def\nman{\numb_{_{\! \manifold}}}
\def\nSph{\numb_{_{\! \Sph}}}
%\def\nSph{\numb_{\operatorname{di}}}
%\def\sph{\omega}
%\def\Sph{\Omega}
\def\sph{\phi}
\def\Sph{\Phi}

\def\hiso{h}
\def\hisom{g}
\def\haniso{b}
\def\aniso{\rho}
\def\anisom{\tau}

\def\BEDR{\cc{B}}
\def\TEDR{\mathbbmsl{T}}
%\def\TEDRm{\TEDR^{\manifold}}
\def\TEDRm{\mathbbmsl{G}}
\def\JEDR{\mathcal{Q}}
\def\REDR{\mathcal{G}}
%\def\REDR{\Gamma}
\def\PrEDR{\Pi}

\def\Kern{\mathcal{K}}
\def\sigmaX{\sigma_{\!_{X}}}
\def\Xbar{\Mv}
\def\corr{\tau}

\def\jm{\!j}
\def\LL{\mathcal{L}}

\def\jd{j^{\circ}}
\def\jjl{j\jl}
\def\jl{\alpha}
\def\weightM{\omega}
\def\svd{\sv^{\circ}}
\def\xvm{\xv^{\dag}}

\def\JJJ{\mathcal{J}}
\def\nJJJ{\numb_{\!_{\JJJ}}}

\def\ini{\operatorname{start}}

%\Section{Efficient dimension reduction in regression}

%%%%%%%%%%%%%%%%%%%%%%%%%%%%%%%  new section  %%%%%%%%%%%%%%%%%%%%%%%%%%%%%%%
\Section{A single-index model}
\label{Ssingind}

Consider the e.d.r. regression problem
\begin{EQA}
	Y_{i}
	&=&
	\fs(\betav^{\T} \Xv_{i}) + \eps_{i} \, ,
\label{YifXieiedrSI}
\end{EQA}
where \( \Xv_{i} \in \R^{\dimd} \) for large \( \dimd \), \( \betav \) is an index vector and
\( \fs \) is a univariate \emph{link function}.
Introduce also a function \( \gs(x) = \E \bigl( Y \cond \Xv=\xv \bigr) \).
A structural single-index assumption yields the relations
\begin{EQA}
	\gs(\xv)
	&=&
	\fs(\betav^{\T} \xv) ,
	\\
	\nabla \gs(\xv)
	&=&
	\fs'(\betav^{\T} \xv) \betav.
\label{fxgKxedrSI}
\end{EQA} 
In particular, the derivative of \( \gs \) at any point \( \xv \) is proportional
to the target vector \( \betav \).
This suggests recovering \( \betav \) from the nonparametrically estimated gradient of \( \gs \).


\Subsection{Preliminaries}
%Define \( \Psi(\xv) = \bigl( \psi_{m}(\xv) \bigr) \in \R^{\dimd+1} \) for \( \psi_{0}(\xv) \equiv 1 \), 
%\(	\psi_{m}(\xv) =	\langle \xv ,\ev_{m} \rangle \), \( m = 1,\ldots,\dimd \). 
For a collection of centers \( \{ \xv_{j} \}_{j \in \JJJ} \) with \( |\JJJ| = \nJJJ \),
localizing tensor \( \TEDR \), 
and a localizing kernel \( \Kern \), consider localizing weights
\begin{EQA}
	\weight_{ij}
	&=&
	\Kern\bigl( \| \TEDR(\Xv_{i} - \xv_{j}) \|^{2} \bigr) 
	\, ,
	\qquad
	i=1,\ldots,n,
	\quad
	j \in \JJJ
	\, .
\label{s7cuqkdwywdhbjwbfjwbdj}
\end{EQA}
At the first step, we apply \( \TEDR = \hiso^{-2} \Id_{\dimp} \)
leading to isotropic weights 
\( \weight_{ij} = \Kern\bigl( \hiso^{-2} \| \Xv_{i} - \xv_{j} \|^{2} \bigr) \).
The bandwidth \( \hiso \) has to be selected to ensure sufficiently many points in almost every vicinity of \( \xv_{j} \)
of radius \( \hiso \). 

Further, for each \( j \in \JJJ \), let \( \Sph_{j} \) be a collection of 
directions (unit vectors) \( \sph \in \R^{\dimd} \) with \( |\Sph_{j}| = \nSph \).
Define for each \( j \in \JJJ \) and \( \sph \in \Sph_{j} \)
\begin{EQA}[rcl]
	\Ima_{j,\sph}
	&=&
	\sumi Y_{i} \langle \Xv_{i} - \xv_{j} \, ,\sph \rangle \, \weight_{ij} 
	\, ,
	\qquad
	S_{j,\sph}
	=
	\sumi \langle \Xv_{i} - \xv_{j} \, ,\sph \rangle \, \weight_{ij}
	\, ,
\label{ytd7e7ydwuhd5eythdY}
%\end{EQ}
%and also the vector \( \Uv_{\jm,\sph} \) in \( \R^{\dimd} \)
%\begin{EQA}[rcl]
	\\
	\Uv_{\jm,\sph}
	&=&
	\sumi (\Xv_{i} - \xv_{j}) \, \langle \Xv_{i} - \xv_{j} \, ,\sph \rangle \, \weight_{ij}
	\in \R^{\dimd}
	\, .
	\qquad
\label{ytd7e7ydwuhd5eythd}
\end{EQA}
%where
%\begin{EQA}[c]
%	N_{j}
%	\eqdef
%	\sumi \weight_{ij}
%	\, ,
%	\quad
%	\Sv_{j}
%	\eqdef
%	\sumi (\Xv_{i} - \xv_{j}) \weight_{ij}
%	\, ,
%	\quad
%	\VP_{j}
%	\eqdef
%	\sumi (\Xv_{i} - \xv_{j})(\Xv_{i} - \xv_{j})^{\T} \weight_{ij}
%	\, .
%\end{EQA}
Model equation \eqref{YifXieiedrSI} implies
\begin{EQA}
	\Ima_{j,\sph}
	&=&
	\sumi Y_{i} \langle \Xv_{i} - \xv_{j} \, ,\sph \rangle \, \weight_{ij}  
	\\
	&=&
	\sumi \fs(\betav^{\T} \Xv_{i}) \langle \Xv_{i} - \xv_{j} \, ,\sph \rangle \, \weight_{ij}
	+ \sumi \eps_{i} \langle \Xv_{i} - \xv_{j} \, ,\sph \rangle \, \weight_{ij} \, .
\label{hgghjkklvddfrdshwey}
\end{EQA}
Further, localization by the weights 
\( \weight_{ij} = \Kern\bigl( \hiso^{-2} \| \Xv_{i} - \xv_{j} \|^{2} \bigr) \) enables us to locally approximate \( \fs(t) \) 
in a vicinity of \( t_{j} = \langle \betav, \xv_{j} \rangle \) by a linear function
\( \fcn_{j} + \fcl_{j} (t - t_{j}) \).
This leads to the approximation 
\begin{EQA}
	&& \nquad
	\sumi  \fs(\betav^{\T} \Xv_{i}) \langle \Xv_{i} - \xv_{j} \, ,\sph \rangle \, \weight_{ij} 
	\approx 
	\sumi \bigl\{ \fcn_{j} + \fcl_{j} (\Xv_{i} - \xv_{j})^{\T} \betav \bigr\} \langle \Xv_{i} - \xv_{j} \, ,\sph \rangle \, \weight_{ij} 
	\\
	&=&
	\fcn_{j} S_{j,\sph} + \fcl_{j} \langle \Uv_{\jm,\sph} \, ,\betav \rangle
%	\KSe_{j} \binom{\fcn_{j}}{\fcl_{j} \, \betav} 
%	=
%	\begin{pmatrix}
%		\fcn_{j} N_{j} + \fcl_{j} \Sv_{j}^{\T} \betav
%		\\
%		\fcn_{j} \Sv_{j} + \fcl_{j} \VP_{j} \betav
%	\end{pmatrix}
	\, .
\label{yd6x5ewbduwuqf64wfgfy6}
\end{EQA}
The calming device suggests introducing additional parameters (observables) 
\( \imav_{j,\sph} = \E \Ima_{j,\sph} \).
The full dimensional parameter vector \( \prmtv \) includes
the target parameter \( \betav \), the observables \( \imav = (\imav_{j,\sph}) \in \R^{\nJJJ \times \nSph} \), 
and the coefficients \( \thetav = (\fcn_{j},\fcl_{j})_{j \in \JJJ} \in \R^{2\nJJJ} \).
The corresponding objective function reads
\begin{EQA}
	\LL(\prmtv)
	&=&
	\LL(\betav,\thetav,\imav)
	=
	\frac{1}{2} \| \Imav - \imav \|^{2} 
	+ \frac{1}{2} \sum_{j \in \JJJ} \sum_{\sph \in \Sph_{j}} 
		\bigl( \ima_{j,\sph} - \fcn_{j} S_{j,\sph} - \fcl_{j} \langle \Uv_{\jm,\sph} \, ,\betav \rangle \bigr)^{2} 
	\\
	&=&
	\frac{1}{2} \sum_{j \in \JJJ} 
	\bigl( \| \Imav_{j} - \imav_{j} \|^{2} + \bigl\| \imav_{j} - \fcn_{j} \Sv_{j} - \fcl_{j} \, \UV_{j} \betav \bigr\|^{2} \bigr) 
	\, ,
\label{gtgyuy6ew5dtde32322jt}
\end{EQA}
where \( \imav_{j} = (\ima_{j,\sph} \, , \sph \in \Sph_{j}) \in \R^{\nSph} \),
\( \Sv_{j} = (S_{j,\sph} \, , \sph \in \Sph_{j}) \in \R^{\nSph} \),
and \( \UV_{j} = (\Uv_{\jm,\sph} \, , \sph \in \Sph_{j})^{\T} \) is a \( \nSph \times \dimd \)-matrix with rows \( \Uv_{\jm,\sph}^{\T} \).
Due to the product \( \fcl_{j} \betav \), this objective function is not quadratic in the scope of parameters.
More precisely, it is a polynomial of order 4 and it can be considered as a special case of the error-in-operator model.
The challenge for a careful analysis is a lack of identifiability: rescaling the vector \( \betav \) by any constant
\( c \) and all the coefficients \( \fcl_{j} \) by \( c^{-1} \) does not change the model equation. 
To overcome this problem and enforce identifiability, we require \( \| \betav \| = 1 \).

\paragraph{Elimination of the observables \( \zv \).}
Representation \eqref{gtgyuy6ew5dtde32322jt} is useful for theoretical analysis.
Numerically, one can eliminate the auxiliary variables \( \zv = (\zv_{j,\sph}) \).

\begin{lemma}
For any \( j \in \JJJ \) and any
\( \fcn_{j},\fcl_{j} \, , \betav \), it holds
\begin{EQA}[c]
	\inf_{\imav_{j} \in \R^{\nSph}} \Bigl( 
		\| \Imav_{j} - \imav_{j} \|^{2} 
		+ \bigl\| \imav_{j} - \fcn_{j} \Sv_{j} - \fcl_{j} \, \UV_{j} \betav \bigr\|^{2} 
	\Bigr) 
	=
	\frac{1}{2} \bigl\| \Imav_{j} - \fcn_{j} \Sv_{j} - \fcl_{j} \, \UV_{j} \betav \bigr\|^{2}
	\, .
\end{EQA}
\end{lemma}
{\color{blue} Please, check.}


This result reduces problem \eqref{gtgyuy6ew5dtde32322jt} to
\begin{EQA}[c]
	\sum_{j \in \JJJ} \bigl\| \Imav_{j} - \fcn_{j} \Sv_{j} - \fcl_{j} \, \UV_{j} \betav \bigr\|^{2}
%	+ \lambda \| \betav - \betav_{0} \|^{2}
	\to
	\inf_{\betav,(\fcn_{j},\fcl_{j})_{j \in \JJJ}}
	\, ,
\label{ikdfjoiefjoieuvbjmr3e}
\end{EQA}
where \( \betav \) is a point on the unit sphere \( \Sphere_{\dimd} \) in \( \R^{\dimd} \).

\paragraph{Relaxation and regularization}
The set \( \Sphere_{\dimd} \) is not convex in \( \R^{\dimd} \), 
and minimization over this set is numerically hard.
We apply the usual relaxation trick by allowing at every step of the minimization procedure 
to consider an unconstrained optimization over \( \betav \in \R^{\dimd} \) and then by projecting 
the obtained solution to the unit sphere.
Additional numerical stability can be ensured 
by a penalty \( \mfrac{\lambda}{2} \| \betav - \betav_{0} \|^{2} \) with \( \betav_{0} \in \R^{\dimd} \)
meaning a previous step solution.
%One can use \( \betav_{0} = (1,0,\ldots,0)^{\T} \) or a random unit vector un \( \R^{\dimd} \).
%It is only important that \( \langle \betav,\ev_{1} \rangle \) significantly deviates from zero.
With a tuning parameter \( \lambda \), this leads to minimization of
\begin{EQA}
	\LL_{\lambda}(\prmtv)
	&=&
	\frac{1}{2} \sum_{j \in \JJJ} \bigl\| \Imav_{j} - \fcn_{j} \Sv_{j} - \fcl_{j} \, \UV_{j} \betav \bigr\|^{2} 
	+ \frac{\lambda}{2} \| \betav - \betav_{0} \|^{2} 
	\, 
\label{tjjc7c6f5fd5e5erjwyhfd}
\end{EQA}
at one step of the procedure.
%

\Subsection{Alternating optimization}
\label{SLLaltern}
Now we describe the alternating optimization procedure for minimization of \( \LL(\prmtv) \).
The idea is to iterate optimization of the objective function w.r.t. \( (\fcn_{j},\fcl_{j})_{j \in \JJJ} \),
for \( \betav \) fixed
and w.r.t. \( \betav \) for \( (\fcn_{j},\fcl_{j})_{j \in \JJJ} \) fixed.
Note that for a fixed \( \betav \), the optimization problem \eqref{ikdfjoiefjoieuvbjmr3e}
can be decomposed into small subproblems 
\begin{EQA}[c]
%	\inf_{(\fcn_{j},\fcl_{j})} 
	\bigl\| \Imav_{j} - \fcn_{j} \Sv_{j} - \fcl_{j} \, \UV_{j} \betav \bigr\|^{2}
	\to \inf_{\fcn_{j},\fcl_{j}}
	\, ,
	\qquad
	j \in \JJJ
	\, .
\end{EQA}
Each of them is quadratic in \( \fcn_{j},\fcl_{j} \) and can be solved analytically.

\begin{lemma}
\label{Lcjlj}
It holds 
\begin{EQA}[c]
	\arginf_{\fcn_{j},\fcl_{j}} 
	\bigl\| \Imav_{j} - \fcn_{j} \Sv_{j} - \fcl_{j} \, \UV_{j} \betav \bigr\|^{2}
	=
	\begin{pmatrix}
		\| \Sv_{j} \|^{2} 
		& 
		\langle \Sv_{j} \, , \UV_{\jm} \betav \rangle
		\\
		\langle \Sv_{j} \, , \UV_{\jm} \betav \rangle  
		& 
		\| \UV_{j} \betav \|^{2} 
	\end{pmatrix}^{-1} 
	\binom{ \langle \Imav_{j} \, , \Sv_{j} \rangle }
	{\langle \Imav_{j} \, , \UV_{\jm} \betav \rangle }
	\, .
\label{ihi88eufiut545yhfee}
\end{EQA}
\end{lemma}
{\color{blue} please check.}

Similarly, for fixed \( \{ (\fcn_{j},\fcl_{j}), j \in \JJJ \} \), minimization of \eqref{ikdfjoiefjoieuvbjmr3e} w.r.t. \( \betav \)
is quadratic and can be solved analytically.

\begin{lemma}
\label{Lbeta}
It holds 
\begin{EQA}[rcl]
	&& \nquad
	\arginf_{\betav \in \R^{\dimd}} \biggl\{ \sum_{j \in \JJJ} \bigl\| \Imav_{j} - \fcn_{j} \Sv_{j} - \fcl_{j} \, \UV_{j} \betav \bigr\|^{2}
	+ \lambda \| \betav - \betav_{0} \|^{2} 
	\biggr\}	
	\\
	&=&
	\biggl\{ \sum_{j \in \JJJ} \fcl_{j}^{2} \, \UV_{j}^{\T} \UV_{j} + \lambda \Id_{\dimd} \biggr\}^{-1}
	\sum_{j \in \JJJ} \bigl\{ \fcl_{j} \, \UV_{j}^{\T} (\Imav_{j} - \fcn_{j} \Sv_{j}) 
	+ \lambda \betav_{0}
	\bigr\}
	\, .
\label{uh8888iiftbe7fuhfff}
\end{EQA}
{\color{blue} please check}
\end{lemma}

\paragraph{Alternating optimization.}
The alternating minimization reads as follows.
\begin{mylist}

\item 
Initialize \( \betav^{(0)} \) randomly or by prior information if available.
Set \( m=1 \).

\item 
With \( \betav^{(m-1)} \) fixed, compute for every \( j \in \JJJ \)
\begin{EQA}[c]
	(\fcn_{j}^{(m)},\fcl_{j}^{(m)})
	=
	\arginf_{\fcn_{j},\fcl_{j}} \bigl\| \Imav_{j} - \fcn_{j} \Sv_{j} - \fcl_{j} \, \UV_{j} \betav^{(m-1)} \bigr\|^{2}
	\, ;
\end{EQA}
see \eqref{ihi88eufiut545yhfee} of Lemma~\ref{Lcjlj}.

\item 
Then with \( (\fcn_{j}^{(m)},\fcl_{j}^{(m)})_{j \in \JJJ} \), update \( \betav^{(m)} \):
\begin{EQA}[c]
	\betav^{(m)}
	=
	\arginf_{\betav} \biggl\{ \sum_{j \in \JJJ} \bigl\| \Imav_{j} - \fcn_{j}^{(m)} \Sv_{j} - \fcl_{j}^{(m)} \, \UV_{j} \betav \bigr\|^{2}
	+ \lambda \| \betav - \betav^{(m-1)} \|^{2} 
	\biggr\}	
	\, ;
\end{EQA}
as suggested in \eqref{uh8888iiftbe7fuhfff} of Lemma~\ref{Lbeta}.

\item
Normalize \( \betav^{(m)} \) to one.
%With \( \rr_{m} \eqdef \| \betav^{(m)} \| \), update \( \betav^{(m)} \to \betav^{(m)}/\rr_{m} \) and \( \fcl_{j} \to \rr_{m} \fcl_{j} \) for
%all \( j \in \JJJ \). 

\item 
Iterate several time for \( m=1,2,\ldots, m_{\max} \).
The procedure delivers the last step unit vector \( \betav^{(m)} \) for \( m = m_{\max} \).
\end{mylist}

{\color{blue} check one-step improvement and convergence of \( \betav^{(m)} \)}

\Subsection{Anisotropic weights and local tensor \( \TEDR \)}
The local weights \( \weight_{ij} \) from \eqref{s7cuqkdwywdhbjwbfjwbdj} involve a local tensor \( \TEDR \).
The single-index model structure \( \gs(\xv) = \fs(\betav^{\T} \xv) \) in \eqref{YifXieiedrSI} suggests
anisotropic weights
\begin{EQA}
	\weights_{ij}
	&=&
	\Kern\bigl( h^{-2} |(\Xv_{i} - \xv_{j})^{\T} \betavs|^{2} \bigr) .
\label{s7cuqkdwywdhbjwbfjwbdji}
\end{EQA}
Unfortunately, such defined weights require knowledge of the index vector \( \betavs \) and 
a correctly specified single-index model. 
We follow the structure-adaptive approach from \cite{HJPS}, which applies anisotropic weights 
with an increasing anisotropy parameter during iterations.
As mentioned earlier, for initialization, we apply a non-informative choice 
\begin{EQA}[c]
	\weight_{ij} = \Kern\bigl( \hiso_{0}^{-2} \| \Xv_{i} - \xv_{j} \|^{2} \bigr) 
	\, .
\end{EQA}
The bandwidth \( \hiso_{0} \) can be selected as the smallest value ensuring
\begin{EQA}[c]
	\frac{1}{\nJJJ}\sum_{j \in \JJJ} \sumi \Kern\Bigl( \mfrac{\| \Xv_{i} - \xv_{j} \|^{2}}{\hiso_{0}^{2}} \Bigr)
%	=
%	\frac{1}{\nJJJ}\sum_{j \in \JJJ} \sumi \weight_{ij}
	\geq 
	\nmin
	\, .
\label{jhgiverkytde4w34etycj}
\end{EQA}
This condition means that in a local vicinity of each point \( \xv_{j} \) there are (in average) 
at least \( \nmin \) design points \( \Xv_{i} \).

Anisotropic weights assume some prior information about the structure of the model, which is accumulated 
in the localizing tensor \( \TEDR \).
Namely, for two values \( \hiso \) and \( \aniso \leq 1 \) and a current guess \( \betav \), anisotropic weights \( \weight_{ij} \)
correspond to the tensor 
\begin{EQA}[rcl]
	\TEDR^{2}
	&=&
	\hiso^{-2} \bigl( \aniso^{2} \, \Id_{\dimd} + \betav \betav^{\T} \bigr)
\label{iufhvu8kjbmnll475edf}
\end{EQA}
leading to 
\begin{EQA}[c]
	\weight_{ij}
	=
	\Kern\bigl( \| \TEDR (\Xv_{i} - \xv_{j}) \|^{2} \bigr)
	=
	\Kern\bigl( 
		\mfrac{\aniso^{2} \| \Xv_{i} - \xv_{j} \|^{2}}{\hiso^{2}} 
		+ \mfrac{\langle \Xv_{i} - \xv_{j}, \betav \rangle^{2}}{\hiso^{2}} 
	\bigr)
	\, .
\end{EQA}
The quantity \( \aniso \) describes the anisotropy strength.
Small values of \( \hiso, \aniso \) correspond to strong localization in direction \( \betav \)
and mild localization in orthogonal directions. 
During iteration, the procedure decreases \( \hiso \) and \( \aniso \) in a way 
to ensure sufficiently many design points an anisotropic local vicinity of \( \xv_{j} \) for in almost every \( j \in \JJJ \).
For fixed \( \hiso \) and \( \betav \), the anisotropy parameter \( \aniso \) is the smallest value ensuring
\begin{EQA}[c]
	\frac{1}{\nJJJ} \sum_{j \in \JJJ} \sumi \Kern\Bigl( 
		\mfrac{\aniso^{2} \| \Xv_{i} - \xv_{j} \|^{2}}{\hiso^{2}} + \mfrac{\langle \Xv_{i} - \xv_{j}, \betav \rangle^{2}}{\hiso^{2}} 
	\Bigr)
%	=
%	\frac{1}{\nJJJ} \sum_{j \in \JJJ} \sumi \weight_{ij}
	\geq 
	\nmin
	\, .
\label{jhgiverkytde4w34etycja}
\end{EQA}
It is easy to see that a decay of \( \hiso \) leads to a decay of \( \aniso \).


\Subsection{Design of experiments and implementation}
Fix the values \( d \) and \( n \) such that \( n \geq 10 d \).
Also, fix the values \( \sigma_{\eps} \), \( \sigmaX \), \( \nJJJ \leq n \), \( \nSph \leq \dimd \), and \( \corr \in [0,1) \).
Finally, fix a unit vector \( \betavs \) and a univariate function \( \fs(x) \); e.g \( f(x) = x^{2} \).
\begin{mylist}

\item 
generate a correlated design \( \Xv_{i} = \sigmaX \bigl\{ \corr \gaussv_{0} + (1 - \corr) \gaussv_{i} \bigr\} \),
where \( \gaussv_{i} \) are i.i.d. standard Gaussian in \( \R^{\dimd} \), \( i=0,1,\ldots,n \);
\item 
generate errors \( \eps_{i} \sim \ND(0,\sigma_{\eps}^{2}) \) and compute \( Y_{i} = \fs(\Xv_{i}^{\T} \betavs) + \eps_{i} \);
\item 
for each \( \Xv_{i} \), decide by Bernoulli experiment with success probability \( \theta \eqdef \nJJJ/n \), whether this point is used for 
constructing the center \( \xv_{j} \) by \( \xv_{j} = \Xv_{i} + \sigmaX \gaussv_{j} \);
ensure \( \#(\xv_{j}) = \nJJJ \);
\item
for each \( j \in \JJJ \), generate the set \( \Sph_{j} \) of unit vectors \( \sph \in \R^{\dimd} \)
with \( |\Sph_{j}| = \nSph \).
One can use random draws \( \sph \eqd \gaussv/\| \gaussv \| \) for \( \gaussv \sim \ND(0,\Id_{\dimd}) \);

%\item
%generate \( \betav_{\ini} \) uniformly on the unit sphere in \( \R^{\dimd} \).
\end{mylist}

\paragraph{Implementation. Tuning parameters.}
For the procedure, we have to fix
\begin{mylist}

%\item
%the value \( \corr \in [0,1) \); par defaut, \( \corr = 1/2 \);
%
\item 
the value \( \nmin \gg 1 \); par defaut, \( \nmin = 10 \);

\item 
the value \( \nJJJ \gg 1 \); par defaut, \( \nJJJ = n \);

\item 
the kernel \( \Kern \); par defaut, Epanechnikov kernel \( \Kern(t) = (1-t^{2})_{+} \);

\item 
the Lagrange penalizing constant \( \lambda \), par defaut, \( \lambda = \nmin \);

\item 
the factor \( a > 1 \); par defaut, \( a = \sqrt{2} \);

\item 
the value \( \hiso_{\min} > 1 \); par defaut, \( \hiso_{\min} = 10 \sigmaX/n \).
\end{mylist}

\paragraph{Implementation. Step 0.}
Fix \( \hiso_{0} \) as the smallest value ensuring
\begin{EQA}[c]
	\frac{1}{\nJJJ} \sum_{j \in \JJJ} 
	\sumi \Kern\Bigl( \mfrac{\| \Xv_{i} - \xv_{j} \|^{2}}{\hiso_{0}^{2}} \Bigr)
%	\bigl( \hiso_{0}^{-2} \| \Xv_{i} - \xv_{j} \|^{2} \bigr)
	\geq 
	\nmin
	\, .
\end{EQA}
%Also, set \( \haniso_{0} = \hiso_{0} \).
\begin{mylist}

\item 
For each \( j \in \JJJ \) and \( \sph \in \Sph_{j} \), 
compute \( \Ima_{j,\sph}^{(0)} \, , S_{j,\sph}^{(0)} \, , \Uv_{\jm,\sph}^{(0)} \) by \eqref{ytd7e7ydwuhd5eythdY} and \eqref{ytd7e7ydwuhd5eythd}
with 
\begin{EQA}[c]
	\weight_{ij} = \Kern\bigl( \hiso_{0}^{-2} \| \Xv_{i} - \xv_{j} \|^{2} \bigr) 
	\, .
\end{EQA}
\item 
%Find \( \prmtv^{(0)} = \{ \betav_{0},(\fcn_{j}^{(0)},\fcl_{j}^{(0)}) \} \) as minimizer of 
%\( \LL_{\lambda}(\prmtv) \) from \eqref{tjjc7c6f5fd5e5erjwyhfd} 
With 
\( \Imav_{j}^{(0)} = \bigl( \Ima_{j,\sph}^{(0)} \bigr)_{\sph \in \Sph_{j}} \),
\( \Sv_{j}^{(0)} = \bigl( S_{j,\sph}^{(0)} \bigr)_{\sph \in \Sph_{j}} \),
\( \UV_{j}^{(0)} = (\Uv_{\jm,\sph}^{(0)} \, , \sph \in \Sph_{j})^{\T} \),
solve
\begin{EQA}[c]
	\bigl\{ \betav_{0},(\fcn_{j}^{(0)},\fcl_{j}^{(0)})_{j \in \JJJ} \bigr\}
	=
	\arginf_{\betav \, , (\fcn_{j},\fclv_{j})_{j \in \JJJ}}
	\sum_{j \in \JJJ} \Bigl( \| \Imav_{j}^{(0)} - \fcn_{j} \Sv_{j}^{(0)} - \fcl_{j} \, \UV_{j}^{(0)} \betav \bigr\|^{2} \Bigr) 
%	+ \frac{\lambda}{2} \| \betav - \betav_{\ini} \|^{2}
	\, 
	\qquad
\end{EQA}
by alternating optimization with a random start \( \betav^{(0)} \); see Section~\ref{SLLaltern}.

%\item 
%%:
%Rescale \( \betav_{0} \) by \( \| \betav_{0} \|^{-1} \) and \( (\fcl_{j}^{(0)}) \) by \( \| \betav_{0} \| \) to ensure 
%\( \| \betav_{0} \| = 1 \).
\end{mylist}

\paragraph{Implementation. Step \( k \).}
For \( k \geq 1 \), 
\begin{mylist}

\item 
update \( \hiso_{k} = \hiso_{k-1} / a \).

\item 
Define \( \TEDR_{k} \) in the form \eqref{iufhvu8kjbmnll475edf} with \( \betav = \betav_{k-1} \) and \( \hiso = \hiso_{k} \):
\begin{EQA}[rcl]
	\TEDR_{k}^{2}
	&=&
	\hiso_{k}^{-2} \bigl( \aniso_{k}^{2} \, \Id_{\dimd} + \betav_{k-1} \betav_{k-1}^{\T} \bigr)
	\, ,
\label{iufhvu8kjbmnll475edfk}
\end{EQA}
where
\( \aniso_{k} \) is fixed by the condition
\begin{EQA}[c]
	\frac{1}{\nJJJ} \sum_{j \in \JJJ} \sumi \Kern\bigl( \| \TEDR_{k} (\Xv_{i} - \xv_{j}) \|^{2}
	\bigr)
%	=
%	\frac{1}{\nJJJ} \sum_{j \in \JJJ} \sumi \Kern\Bigl( 
%		\mfrac{\aniso_{k}^{2} \| \Xv_{i} - \xv_{j} \|^{2}}{\hiso_{k}^{2}} 
%		+ \mfrac{\langle \Xv_{i} - \xv_{j}, \betav_{k-1} \rangle^{2}}{\hiso_{k}^{2}} 
%	\Bigr)
	\geq 
	\nmin
	\, .
	\qquad 
\label{jhgiverkytde4w34etycja}
\end{EQA}
{\color{blue} plot \( \aniso_{k} \)}

\item
For each \( j \in \JJJ \),
update the sets of directions \( \Sph_{j} \) by drawing \( \sph \eqd \gaussv_{\TEDR_{k}}/\| \gaussv_{\TEDR_{k}} \| \)
for \( \gaussv_{\TEDR_{k}} \sim \ND(0,\TEDR_{k}^{2}) \);

\item 
For each \( j \in \JJJ \), \( \sph \in \Sph_{j} \), 
compute \( \Ima_{j,\sph}^{(k)} \, , S_{j,\sph}^{(k)} \, , \Uv_{\jm,\sph}^{(k)} \)
by \eqref{ytd7e7ydwuhd5eythdY} and \eqref{ytd7e7ydwuhd5eythd} with
\begin{EQA}[c]
	\weight_{ij} 
	= 
	\weight_{ij}^{(k)}
	= 
	\Kern\bigl( \| \TEDR_{k} (\Xv_{i} - \xv_{j}) \|^{2} \bigr) 
	\, .
\end{EQA}

\item 
With 
\( \Imav_{j}^{(k)} = \bigl( \Ima_{j,\sph}^{(k)} \bigr)_{\sph \in \Sph_{j}} \),
\( \Sv_{j}^{(k)} = \bigl( S_{j,\sph}^{(k)} \bigr)_{\sph \in \Sph_{j}} \),
\( \UV_{j}^{(k)} = (\Uv_{\jm,\sph}^{(k)} \, , \sph \in \Sph_{j})^{\T} \),
solve
\begin{EQA}[c]
	\bigl\{ \betav_{k},(\fcn_{j}^{(k)},\fcl_{j}^{(k)})_{j \in \JJJ} \bigr\}
	=
	\arginf_{\betav \, , (\fcn_{j},\fclv_{j})_{j \in \JJJ}}
	\sum_{j \in \JJJ} \Bigl( \| \Imav_{j}^{(k)} - \fcn_{j} \Sv_{j}^{(k)} - \fcl_{j} \, \UV_{j}^{(k)} \betav \bigr\|^{2} \Bigr) 
%	+ \frac{\lambda}{2} \| \betav - \betav_{k-1} \|^{2}
	\, ,
	\qquad
\end{EQA}
using alternating optimization from Section~\ref{SLLaltern} which starts at \( \betav^{(0)} = \betav_{k-1} \).
%\item 
%%:
%Rescale \( \betav_{k} \) by \( \| \betav_{k} \|^{-1} \) and \( (\fcl_{j}^{(k)}) \) by \( \| \betav_{k} \| \) to ensure 
%\( \| \betav_{k} \| = 1 \).


\item
If \( \hiso_{k}/a > \hiso_{\min} \) then increase \( k \) by one and loop.

\item
Otherwise terminate with output \( \betav^{(k)} \).
% and \( (\fcn_{j}^{(k)},\fcl_{j}^{(k)})_{j \in \JJJ} \).
\end{mylist}


{\color{blue} 
\begin{mylist}
\item
Fix \( \dimd = 10 \) and explore the performance of the method 
for different values of tuning parameters, especially \( \nmin \), \( \nSph \), and \( \lambda \);

\item
behavior for different \( \sigma_{\eps} \), \( \sigmaX \), \( q \in (0,1] \), and \( \corr \in [0,1) \);

\item
try with different \( \dimd \) and find a breakdown dimension \( \dimd \).

\end{mylist}
 
}



%%%%%%%%%%%%%%%%%%%%%%%%%%%%%%%  new section  %%%%%%%%%%%%%%%%%%%%%%%%%%%%%%%
\Section{Multi-index model}
\label{Smultind}


Consider the e.d.r. regression problem
\begin{EQA}
	Y_{i}
	&=&
	\fs(\PEDR \Xv_{i}) + \eps_{i}
	\, ,
\label{YifXieiedr}
\end{EQA}
where \( \Xv_{i} \in \R^{\dimd} \) for large \( \dimd \), 
\( \PEDR \) is an unknown element from the set \( \PROJ_{\dimd,\dimm} \) 
of orthogonal mappings from 
from \( \R^{\dimd} \) to \( \R^{\dimm} \) for \( \dimm \ll \dimd \), and
\( \fs \) is a function on \( \R^{\dimm} \).
The mapping \( \PEDR \) can be written as 
\begin{EQA}
	\PEDR \xv
	&=&
	(\betav_{1}^{\T} \xv) \betav_{1} + \ldots + (\betav_{\dimm}^{\T} \xv) \betav_{\dimm}
	\, ,
\label{ydgenb8jejohkou90tkr}
\end{EQA}
where \( \betav_{m}^{\T} \) are the orthonormal rows of \( \PEDR \).
Here the identifiability issue is even more severe than in the single-index model, because 
the set of the vectors \( \betav_{j} \) can be exchanged, rotated, etc. 

Similarly to the single-index case, 
the structural information about the model
can be recoved from the gradient vectors \( \nabla \fs(\xv) \).
More precisely, under \eqref{ucjcvjuvyvyf6r6rwjdy}, the matrix
%by orthogonality of \( \betav_{1}, \ldots,\betav_{m} \)
\begin{EQA}[c]
	\JEDR
	\eqdef
	\sum_{j \in \JJJ} \nabla \fs(\xv_{j}) \, \nabla \fs(\xv_{j})^{\T}
	=
	\sum_{j \in \JJJ} \PEDR^{\T} \fclv_{j} \, \bigl( \PEDR^{\T} \fclv_{j} \bigr)^{\T}
%	\biggl( \sum_{m=1}^{\dimm} \fcl_{j,m} \, \betav_{m} \biggr)^{\T}
	=
	\PEDR^{\T} \biggl( \sum_{j \in \JJJ} \fclv_{j} \, \fclv_{j}^{\T} \biggr) \PEDR
	\, 
\label{ihjfiojihfer6t4rrt}
\end{EQA}
describes the multi-index structure.
Moreover, 
If all positive eigenvalues of \( \JEDR \) are different, 
then the SVD \( \JEDR = \PEDR^{\T} \Lambda \, \PEDR \) with \( \Lambda = \diag(\lambda_{1} \geq \ldots \geq \lambda_{\dimm}) \)
provides a unique canonical representation for the multi-index model \eqref{YifXieiedr}.


\Subsection{Preliminaries}
For a collection of centers \( \{ \xv_{j} \}_{j \in \JJJ} \),
and a localizing kernel \( \Kern \), consider the localizing weights of the form
\begin{EQA}
	\weight_{ij}
	&=&
	\Kern\bigl( \| \TEDR(\Xv_{i} - \xv_{j}) \|^{2} \bigr) 
	\, ,
\label{s7cuqkdwywdhbjwbfjwbdjm}
\end{EQA}
where \( \TEDR \) is a localizing tensor \( \TEDR \) (symmetric \( \dimd \)-matrix)
which reflects the available information about the multi-index structure in \eqref{YifXieiedr}.
At the beginning, an uninformative isotropic tensor \( \TEDR_{0} = \hiso_{0}^{-2} \Id_{\dimd} \) will be used.
During iteration, we learn the structure and use it to improve \( \TEDR \); see Section~\ref{SEDRsa} ahead.
Further, for every \( j \in \JJJ \), assume a set \( \Sph_{j} \) or directions \( \sph \) (unit vectors in \( \R^{\dimd} \)) 
to be fixed.
Compute for all \( j \in \JJJ \), \( \sph \in \Sph_{j} \),
\begin{EQ}[rcccl]
	\Ima_{j,\sph}
	&=&
	\sumi Y_{i} \langle \Xv_{i} - \xv_{j} \, ,\sph \rangle \, \weight_{ij} 
	\, ,
	&
	\\
	S_{j,\sph}
	&=&
	\sumi \langle \Xv_{i} - \xv_{j} \, ,\sph \rangle \, \weight_{ij}
	\, ,
	&
%\label{ytd7e7ydwuhd5eythdYm}
	\quad
	\Uv_{\jm,\sph}
	&=
	\sumi (\Xv_{i} - \xv_{j}) \, \langle \Xv_{i} - \xv_{j} \, ,\sph \rangle \, \weight_{ij}
	\, .
	\qquad
\label{ytd7e7ydwuhd5eythdm}
\end{EQ}
Model equation \eqref{YifXieiedrSI} implies
\begin{EQA}
	\Ima_{j,\sph}
	&=&
	\sumi Y_{i} \langle \Xv_{i} - \xv_{j} \, ,\sph \rangle \, \weight_{ij}  
	\\
	&=&
	\sumi \fs(\PEDR \Xv_{i}) \langle \Xv_{i} - \xv_{j} \, ,\sph \rangle \, \weight_{ij}
	+ \sumi \eps_{i} \langle \Xv_{i} - \xv_{j} \, ,\sph \rangle \, \weight_{ij} \, .
\label{hgghjkklvddfrdshweym}
\end{EQA}
%
%as in \eqref{ytd7e7ydwuhd5eythd}.
The function \( \fs \) is now \( \dimm \)-variate, and we consider a linear approximation 
\begin{EQA}
	\fs(\PEDR \xv) 
	& \approx &
	\fcn_{j} + \fclv_{j}^{\T} \PEDR (\xv - \xv_{j})
	=
	\fcn_{j} + \biggl( \sum_{m=1}^{\dimm} \fcl_{j,m} \, \betav_{m} \biggr)^{\T} (\xv - \xv_{j}) ,
\label{ucjcvjuvyvyf6r6rwjdy}
\end{EQA}
where \( \fclv_{j} = (\fcl_{j,m}) \in \R^{\dimm} \).
Correspondingly
\begin{EQA}[rcl]
	&& \nquad
	\sumi \fs(\PEDR \Xv_{i}) \, \langle \Xv_{i} - \xv_{j} \, ,\sph \rangle \, \weight_{ij} 
	\approx 
	\sumi \bigl\{ \fcn_{j} 
%	+ \sum_{m=1}^{\dimm} \fcl_{j,m} \, (\Xv_{i} - \xv_{j})^{\T} \betav_{m} \Bigr\}
	+ \fclv_{j}^{\T} \PEDR (\Xv_{i} - \xv_{j})
	\langle \Xv_{i} - \xv_{j} \, ,\sph \rangle \bigr\} \weight_{ij} 
	\\
	&=&
	\fcn_{j} S_{j,\sph} + \fclv_{j}^{\T} \PEDR \, \Uv_{\jm,\sph}
	=
	\fcn_{j} S_{j,\sph} + \Uv_{\jm,\sph}^{\T} \PEDR^{\T} \, \fclv_{j}
	\, .
\label{yd6x5ewbduwuqf64wfgfymi}
\end{EQA}
%where \( \fclv_{j} \cdot \betav = \fcl_{j,1} \, \betav_{1} + \ldots + \fcl_{j,\dimm} \, \betav_{\dimm} \).
The full dimensional parameter \( \prmtv \) includes
%the corresponding loading factors \( \mue_{1},\ldots,\mue_{\dimm} \), 
\( (\imav_{j,\sph}) \in \R^{\nJJJ \nSph} \), 
the coefficients \( (\fcn_{j},\fclv_{j})_{j \in \JJJ} \in \R^{\nJJJ(\dimm+1)} \),
and the target mapping \( \PEDR = (\betav_{1},\ldots,\betav_{\dimm})^{\T} \).
Denote \( \imav_{j} = \bigl( \ima_{j,\sph} \bigr)_{\sph \in \Sph_{j}} \in \R^{\nSph} \), 
\( \Imav_{j} = \bigl( \Ima_{j,\sph} \bigr)_{\sph \in \Sph_{j}} \in \R^{\nSph} \),
\( \Sv_{j} = \bigl( S_{j,\sph} \bigr)_{\sph \in \Sph_{j}} \in \R^{\nSph} \),
and let \( \UV_{j} \) be the \( \nSph \times \dimd \)-matrix with the rows \( \Uv_{\jm,\sph}^{\T} \),
\( \sph \in \Sph_{j} \).
Then the objective function reads
\begin{EQA}
	\LL(\prmtv)
	&=&
	\frac{1}{2} \sum_{j \in \JJJ} \Bigl( \| \Imav_{j} - \imav_{j} \|^{2} 
		+ \bigl\| \imav_{j} - \fcn_{j} \Sv_{j} - \UV_{j} \, \PEDR^{\T} \fclv_{j} \bigr\|^{2} \Bigr) 
	\, .
	\qquad
%\label{gtgyuy6ew5dtde32322jtmi}
\end{EQA}
The auxillary variables \( \zv_{j} \) can be eliminated, leading to
\begin{EQA}[rcl]
	\LL(\prmtv)
	&=&
	\LL(\PEDR,(\fcn_{j},\fclv_{j})_{j \in \JJJ})
	=
	\frac{1}{4} \sum_{j \in \JJJ} \bigl\| \Imav_{j} - \fcn_{j} \Sv_{j} - \UV_{j} \, \PEDR^{\T} \fclv_{j} \bigr\|^{2}	
	\, .
	\qquad
\label{gtgyuy6ew5dtde32322jtmiG0}	
\end{EQA}
This function should be minimized over the set of all coefficients \( (\fcn_{j},\fclv_{j})_{j \in \JJJ} \) and all
orthonormal mappings \( \PEDR \in \PROJ_{\dimd,\dimm} \) with \( \PEDR \PEDR^{\T} = \Id_{\dimm} \):
\begin{EQA}[c]
%	\prmtv_{0}
%	= 
	\{ \PEDR_{0} \, , (\fcn_{j}^{(0)},\fclv_{j}^{(0)})_{j \in \JJJ} \} 
	=
	\arginf_{\PEDR \, , (\fcn_{j},\fclv_{j})_{j \in \JJJ}} \LL(\prmtv)
	\, .
\label{ju7fudeju74334fhfd}
\end{EQA}
The class of mappings \( \PROJ_{\dimd,\dimm} \) is not convex, this makes the problem computationally hard.
We also face an identifiability issue: a solution \( \PEDR_{0} \) corresponds to 
a particular choice of the basis in the EDR space \( \EDR \).
This basis can be arbitrarily rotated.

\Subsection{Alternating optimization}
\label{SaltMI}
For solving \eqref{ju7fudeju74334fhfd}, one can apply the alternating procedure similar to Section~\ref{SLLaltern}.
The basic idea of the proposed approach is to relax the target \( \PEDR \)
and optimize \( \LL(\prmtv) \) w.r.t. 
\( \BEDR = (\betav_{1},\ldots,\betav_{\dimm})^{\T} \in \R^{\dimm \times \dimd} \). 
Then the desired orthogonal mapping \( \PEDR \) will be obtained using SVD 
of the solution \( \BEDR \).
%
To ensure numerical stability, one can add a penalty term \( \lambda \| \BEDR - \PEDR_{m-1} \|_{\Fr}^{2} \)
for a previous step solution \( \PEDR_{m-1} \).
%This leads to
%\begin{EQA}
%	\LL_{\lambda}(\prmtv)
%	&=&
%	\frac{1}{4} \sum_{j \in \JJJ} \bigl\| \Imav_{j} - \fcn_{j} \Sv_{j} - \UV_{j} \, \BEDR^{\T} \fclv_{j} \bigr\|^{2} 
%	+ \frac{\lambda}{2} \| \BEDR - \PEDR_{0} \|_{\Fr}^{2}
%	\, ,
%	\qquad
%\label{gtgyuy6ew5dtde32322jtmiG}
%\end{EQA}
%which has to be minimized over \( \BEDR \in \R^{\dimd \times \dimm} \) and 
%\( (\fcn_{j},\fclv_{j})_{j \in \JJJ} \in \R^{\dimm+1} \).

The alternating minimization reads as follows.
\paragraph{Initialization}
Set \( m=0 \) and fix
\( \PEDR_{0} \) randomly or using a previous step estimate.

\paragraph{Loop}
\begin{mylist}

\item 
Increase \( m \) by one. 

\item
With \( \PEDR_{m-1} \) fixed, compute for every \( j \in \JJJ \)
\begin{EQA}[c]
	(\fcn_{j}^{(m)},\fclv_{j}^{(m)})
	=
	\arginf_{\fcn_{j},\fclv_{j}} \bigl\| \Imav_{j} - \fcn_{j} \Sv_{j} - \UV_{j} \PEDR_{m-1}^{\T} \fclv_{j} \bigr\|^{2}
\end{EQA}
similarly to Lemma~\ref{Lcjlj}.
{\color{blue} please provide a closed form}

\item 
Then, with \( (\fcn_{j}^{(m)},\fclv_{j}^{(m)})_{j \in \JJJ} \) fixed, compute \( \BEDR_{m} \):
\begin{EQA}[c]
	\BEDR_{m}
	=
	\arginf_{\BEDR \in \R^{\dimd \times \dimm}} 
		\biggl\{ \sum_{j \in \JJJ} \bigl\| \Imav_{j} - \fcn_{j}^{(m)} \Sv_{j} - \UV_{j} \BEDR^{\T} \fclv_{j}^{(m)} \bigr\|^{2}
	+ \lambda \| \BEDR - \PEDR_{m-1} \|^{2} 
	\biggr\}	
	\, ,
\end{EQA}
similarly to Lemma~\ref{Lbeta}. {\color{blue} please provide a closed form}

\item
Compute 
\begin{EQA}[rcl]
	\JEDR_{m}
	\eqdef
	\sum_{j \in \JJJ} (\BEDR_{m}^{\T} \, \fclv_{j}^{(m)}) \, \bigl( \BEDR_{m}^{\T} \, \fclv_{j}^{(m)} \bigr)^{\T}
	&=&
	\BEDR_{m}^{\T} \Bigl( \sum_{j \in \JJJ} {\fclv}_{j}^{(m)} \, ({\fclv}_{j}^{(m)})^{\T} \Bigr) \BEDR_{m}
\end{EQA}
and its SVD
\begin{EQA}[c]
	\JEDR_{m}
	=
	\PEDR_{m}^{\T} \, \Lambda_{m} \, \PEDR_{m}
	\, ,
	\qquad
	\PEDR_{m} \, \PEDR_{m}^{\T}
	= 
	\Id_{\dimm}
	\, ,
	\qquad
	\Lambda_{m} = \diag(\lambda_{1} \geq \ldots \geq \lambda_{\dimm})
	\, .
\end{EQA}

\item 
Iterate several times until \( m=m_{\max} \).
Output \( \PEDR_{m} \) and \( \JEDR_{m} \) for \( m = m_{\max} \).
\end{mylist}

{\color{blue} check one-step improvement and convergence of \( \PEDR_{m} \)}



\Subsection{Structural adaptation}
\label{SEDRsa}
The use of isotropic weights \( \weight_{ij} \) for computing the observables 
\( \Imav_{j} \), \( \Sv_{j} \), and \( \UV_{j} \) is suboptimal and leads to poor 
quality of recovering the EDR-space \( \EDR \).
The structure-adaptive approach from \cite{HJPS} suggests to extract the structural information from the current-step estimate 
and use for improving the accuracy of the next-step estimate by using 
anisotropic weights \( \weight_{ij} \).

Alternating procedure from Section~\ref{SaltMI} delivers 
a starting solution \( \PEDR_{0} \) of \eqref{ju7fudeju74334fhfd} and also the estimate \( \JEDR_{0} \) of \( \JEDR \)
from \eqref{ihjfiojihfer6t4rrt}.
These matrices can be used to improve the procedure by redefining the tensor \( \TEDR \) as
\begin{EQA}
	\TEDR^{2}
	&=&
	\hiso^{-2} \bigl( \aniso^{2} \, \Id_{\dimd} + \JEDR_{0} \bigr)
%	\sum_{m} \biggl( \sum_{j \in \JJJ} \fcl_{j,m}^{2} \biggr) \betav_{m} \betav_{m}^{\T} 
	\, ,
\label{yfijowjdo76dvuienk}
\end{EQA}
with \( \aniso < 1 \).
This value describes the degree of kernel anisotropy.
Alternatively, %one can use
\begin{EQA}[c]
	\TEDR^{2}
	=
	\hiso^{-2} \aniso^{2} (\Id_{\dimd} - \PEDR_{0}^{\T} \PEDR_{0}) + \hiso^{-2} \, \JEDR_{0}
	\, .
\end{EQA}
This tensor leads to anisotropic weights 
\( \weight_{ij} = \Kern\bigl( \| \TEDR(\Xv_{i} - \xv_{j}) \|^{2} \bigr) \)
which strongly penalize in the principle directions of the estimated EDR space.
It will also be used for reneuval of the set of directions \( \Sph_{j} \) which are randomly sampled as
\( \gaussv_{\TEDR}/\| \gaussv_{\TEDR} \| \) for \( \gaussv_{\TEDR} \sim \ND(0,\TEDR^{2}) \). 
The sums \( \Ima_{j,\sph} \), \( S_{j,\sph} \), and \( \Uv_{\jm,\sph} \) in \eqref{ytd7e7ydwuhd5eythdm}
have to be recomputed with these changes.
The improved estimator \( \tilde{\prmtv} \) is given by 
minimization of \( \LL(\prmtv) \) from \eqref{gtgyuy6ew5dtde32322jtmiG0}
after this update.
%
This structural adaptation step can be iterated several times.
The values \( \hiso, \aniso \) may change during iterations: 
\( \hiso \) decreases at every step with a fixed factor 
\( a = 2^{1/\dimm} \), while \( \aniso \) is fixed by condition similar to \eqref{jhgiverkytde4w34etycja}.


\Subsection{Implementation}
\label{SEDRmiproc}
Fix 
\( \dimm (\leq 3) \), 
\( \nmin (= 10) \), 
\( \nJJJ (= n) \), 
\( \nSph (= \nmin \, \dimm) \),
\( a (= 2^{1/\dimm}) \), 
\( \lambda (= \dimm) \),
\( m_{\max} (= 5) \).

\paragraph{Step 0}
\begin{mylist}
\item
Generate \( (\xv_{j} \, , j \in \JJJ) \).

\item
Fix \( \hiso_{0} \) as the smallest value ensuring
\begin{EQA}[c]
	\frac{1}{\nJJJ} \sum_{j \in \JJJ} 
	\sumi \Kern\bigl( \hiso_{0}^{-2} \| \Xv_{i} - \xv_{j} \|^{2} \bigr)
%	=
%	\frac{1}{\nJJJ} \sum_{j \in \JJJ} \sumi \weight_{ij}
	\geq 
	\nmin
	\, .
\end{EQA}

\item
For each \( j \in \JJJ \), generate the set \( \Sph_{j} \) of \( \nSph \) unit vectors \( \sph \) in \( \R^{\dimd} \).

\item 
For each \( j \in \JJJ \) and \( \sph \in \Sph_{j} \), compute  
\( \Ima_{j,\sph}^{(0)} \, , \, S_{j,\sph}^{(0)} \, , \, \Uv_{\jm,\sph}^{(0)} \)
by \eqref{ytd7e7ydwuhd5eythdm} with 
\begin{EQA}[c]
	\weight_{ij} 
	= 
	\weight_{ij}^{(0)}
	= 
	\Kern\bigl( \hiso_{0}^{-2} \| \Xv_{i} - \xv_{j} \|^{2} \bigr) 
	\, .
\end{EQA}

\item 
For
\( \Imav_{j}^{(0)} = \bigl( \Ima_{j,\sph}^{(0)} \bigr)_{\sph \in \Sph_{j}} \),
\( \Sv_{j}^{(0)} = \bigl( S_{j,\sph}^{(0)} \bigr)_{\sph \in \Sph_{j}} \),
\( \UV_{j}^{(0)} = (\Uv_{\jm,\sph}^{(0)} \, , \sph \in \Sph_{j})^{\T} \),
solve 
%\( \Imav_{j}^{(0)} = (\Ima_{j,\sph}^{(0)})_{\sph \in \Sph_{j}} \),
%\( \Sv_{j}^{(0)} = (S_{j,\sph}^{(0)})_{\sph \in \Sph_{j}} \),
%\( \UV_{j}^{(0)} = (\Uv_{\jm,\sph}^{(0)})^{\T} \), find 
\begin{EQA}[c]
	\{ \PEDR_{0} \, , (\fcn_{j}^{(0)},\fclv_{j}^{(0)})_{j \in \JJJ} \} 
	= 
	\arginf_{\PEDR \, , (\fcn_{j},\fclv_{j})_{j \in \JJJ}} 
	\sum_{j \in \JJJ} \bigl\| \Imav_{j}^{(0)} - \fcn_{j} \Sv_{j}^{(0)} - \UV_{j}^{(0)} \, \PEDR^{\T} \fclv_{j} \bigr\|^{2}
	\, 
	\qquad
\label{jhvihjikr545yhvcvv}
\end{EQA}
by alternating optimization from Section~\ref{SaltMI} with random initialization of \( \PEDR^{(0)} \). 
%Here, one can use a random initialization of \( \PEDR \).
%The procedure delivers \( \PEDR_{0} \) and \( \JEDR_{0} \) which will be used for structural adaptation.
%The solution is not unique; any minimizer can be used.  
%from \eqref{gtgyuy6ew5dtde32322jtmi by alternating minimization.

%\item 
%Compute
%\begin{EQA}[rcl]
%	\JEDR_{0}
%	\eqdef
%	\sum_{j \in \JJJ} (\BEDR_{0}^{\T} {\fclv}_{j}^{(0)}) \, \bigl( \BEDR_{0}^{\T} {\fclv}_{j}^{(0)} \bigr)^{\T}
%	&=&
%	\BEDR_{0}^{\T} \Bigl( \sum_{j \in \JJJ} {\fclv}_{j}^{(0)} \, ({\fclv}_{j}^{(0)})^{\T} \Bigr) \BEDR_{0}
%%	\, .
%\end{EQA}
%and its SVD
%\begin{EQA}[c]
%	\JEDR_{0}
%	=
%	\PEDR_{0}^{\T} \Lambda_{0} \, \PEDR_{0}
%	\, ,
%	\qquad
%	\PEDR_{0} \, \PEDR_{0}^{\T}
%	= 
%	\Id_{\dimm}
%	\, ,
%	\quad
%	\Lambda_{0} = \diag(\lambda_{1} \geq \ldots \geq \lambda_{\dimm})
%	\, .
%\end{EQA}
\item Set \( k=1 \).
\end{mylist}


\paragraph{Step \( k \), assumed \( \PEDR_{k-1} \) and \( \JEDR_{k-1} \) available.}

\begin{mylist}

\item
Update \( \hiso_{k} = \hiso_{k-1} / a \).

\item
Define \( \TEDR_{k} \) by
\begin{EQA}[c]
	\TEDR_{k}^{2}
	\eqdef
	\hiso_{k}^{-2} \aniso_{k}^{2} \, \Id_{\dimd} + \hiso_{k}^{-2} \, \JEDR_{k-1}
	\qquad
	\text{or }
	\quad
	\hiso_{k}^{-2} \aniso_{k}^{2} (\Id_{\dimd} - \PEDR_{k-1}^{\T} \PEDR_{k-1}) + \hiso_{k}^{-2} \, \JEDR_{k-1}
	\, ,
\end{EQA}
{\color{blue} please, implement and compare both options}
\\
where \( \aniso_{k} \) is the smallest value ensuring 
\begin{EQA}[c]
	\frac{1}{\nJJJ} \sum_{j \in \JJJ} \sumi 
		\Kern\bigl( \| \TEDR_{k} (\Xv_{i} - \xv_{j}) \|^{2} \bigr)
%	=
%	\frac{1}{\nJJJ}\sum_{j \in \JJJ} \sumi \weight_{ij}
	\geq 
	\nmin
	\, .
\label{jhgiverkytde4w34etycjam}
\end{EQA}

\item
For each \( j \in \JJJ \), 
update the sets of directions \( \Sph_{j} \) by drawing \( \sph \eqd \gaussv_{\TEDR_{k}}/\| \gaussv_{\TEDR_{k}} \| \)
for \( \gaussv_{\TEDR_{k}} \sim \ND(0,\TEDR_{k}^{2}) \).

\item 
Compute \( \bigl( \Ima_{j,\sph}^{(k)} \, , \, S_{j,\sph}^{(k)} \, , \, \Uv_{\jm,\sph}^{(k)} \bigr)_{j \in \JJJ , \, \sph \in \Sph_{j}} \) 
by \eqref{ytd7e7ydwuhd5eythdm} with \( \weight_{ij} = \Kern\bigl( \| \TEDR_{k} (\Xv_{i} - \xv_{j}) \|^{2} \bigr) \).

\item
With 
\( \Imav_{j}^{(k)} = \bigl( \Ima_{j,\sph}^{(k)} \bigr)_{\sph \in \Sph_{j}} \),
\( \Sv_{j}^{(k)} = \bigl( S_{j,\sph}^{(k)} \bigr)_{\sph \in \Sph_{j}} \),
\( \UV_{j}^{(k)} = (\Uv_{\jm,\sph}^{(k)} \, , \sph \in \Sph_{j})^{\T} \),
solve
\begin{EQA}[c]
	(\PEDR_{k},(\fcn_{j}^{(k)},\fclv_{j}^{(k)})_{j \in \JJJ})
	=
	\arginf_{\PEDR,(\fcn_{j},\fclv_{j})_{j \in \JJJ}}
	\sum_{j \in \JJJ} \bigl\| \Imav_{j}^{(k)} - \fcn_{j} \Sv_{j}^{(k)} - \UV_{j}^{(k)} \, \PEDR^{\T} \fclv_{j} \bigr\|^{2} 
%	+ \lambda \| \BEDR - \PEDR_{k-1} \|_{\Fr}^{2}
	\, 
\end{EQA}
%where \( \PEDR %= (\betav_{1}, \ldots, \betav_{\dimm})^{\T} \in \R^{\dimm \times \dimd} \), 
%\in \PROJ_{\dimd, \dimm} \), 
%\( \fcn_{j} \in \R \), \( \fclv_{j} \in \R^{\dimm} \), \( j \in \JJJ \),
using the alternating optimization as in Section~\ref{SaltMI} starting from \( \PEDR_{k-1} \).

\item
Terminate if \( \hiso_{k}/a < \hiso_{\min} \).
Output \( \PEDR_{k} \), \( \JEDR_{k} \).
%, \( (\fcn_{j}^{(k)},\fclv_{j}^{(k)})_{j \in \JJJ} \).

\item 
%:
Otherwise, increase \( k \) by one and loop. 

\end{mylist}

{\color{blue} 
\begin{mylist}
\item
The first step is important. 
If \( \cos(\PEDRs,\PEDR_{0}) \approx 0 \) then no sense to continue;
try with new directions \( \Sph_{j} \).

\item
Parameters \( \nSph \), \( \lambda \), are especially questionable;

\item
check whether renewal of directions \( \sph \in \Sph_{j} \) is helpful.


\end{mylist}

}

%%%%%%%%%%%%%%%%%%%%%%%%%%%%%%%  new section  %%%%%%%%%%%%%%%%%%%%%%%%%%%%%%%
\Section{Manifold learning}

Consider a regression \( Y_{i} = \fs(\Xv_{i}) + \eps_{i} \) with a high-dimensional \( \Xv_{i} \in \R^{\dimd} \),
\( n = 1,\ldots,n \).
The manifold hypothesis assumes that the function \( \fs(\cdot) \) has a nearly multi-index structure in a local vicinity 
of each point \( \xv \in \R^{\dimd} \).
This means a local linear approximation
\begin{EQA}[c]
	\E Y_{i}
	\approx
	\fcn_{j} + \fclv_{j}^{\T} \PEDR_{j} (\Xv_{i} - \xv_{j})
\end{EQA}
for all points \( \Xv_{i} \) in a vicinity of \( \xv_{j} \).
Here not only the intercept \( \fcn_{j} \in \R \) and 
the slope coefficients \( \fclv_{j} = (\fcl_{1},\ldots,\fcl_{\dimm}) \in \R^{\dimm} \)
but also the mapping \( \PEDR_{j} = \PEDR(\xv) \in \PROJ_{\dimd,\dimm} \)
depend on \( \xv_{j} \).
%
The proposed approach extends the multi-index procedure from Section~\ref{Smultind} to this situation 
by considering a family \( (\PEDR_{j}) \) in place of one mapping \( \PEDR \) with a strong penalty on their variability.
This approach does not require any manifold description.
We also do not project the data on the manifold.

\Subsection{Multiscale modeling}
The manifold hypothesis means that for any fixed \( \xvd \) and any \( \xv \) close to \( \xvd \), 
one can use the mapping \( \PEDR(\xvd) \) in place of \( \PEDR(\xv) \). 
%
This discussion suggests to introduce two scales.
One of them is for local linear approximation of the function \( \fs(\cdot) \) in the form
\begin{EQA}[c]
	\fs(\xv) - \fs(\xvd)
	\approx
	\fclv^{\T} \, \PEDR (\xv - \xvd) 
\end{EQA} 
with the coefficients \( \fclv = (\fcl_{m}) \) and mappings \( \PEDR \in \PROJ_{\dimd,\dimm} \) depending
on the central point \( \xvd \).
Another one for local manifold approximation:  
\(
	\PEDR(\xv)
	\approx
	\PEDR(\xvd)
%	\, .
\).
%A local tensor 
%\( \TEDR_{j} \) and bandwidth \( \hiso_{0} \) for the function \( \fs \),
%and tensor \( \TEDRm_{j} \) and bandwidth \( \hisom \) for the manifold.
At the injitialization step of the procedure, we consider isotropic vicinity
\( \{ \| \xv - \xvd \| \leq \hiso_{0} \} \) for approximation of \( \fs \)
and \( \{ \| \xv - \xvd \| \leq \hisom_{0} \} \) for continuity of \( \PEDR(\xv) \).
The approach assumes that the manifold is ``smoother'' than the function \( \fs \) and hence,
\( \hisom_{0} \gg \hiso_{0} \).
%, where \( \hiso_{0} \) is used for the weights \( \weight_{ij} \).

\Subsection{Objective function and manifold penalty}
The approach extends the procedure for the multi-index case from Section~\ref{SEDRmiproc}.
The basic idea behind the procedure from \eqref{ju7fudeju74334fhfd} is that the EDR space \( \EDR \)
is the same for all \( \xv_{j} \).
This structural assumption does not hold in the manifold case.
However, it is nearly fulfilled if we limit the analysis to a vicinity around each \( \xv_{j} \).
This suggests using two hierarchic localizing schemes described by 
localizing tensors \( \TEDR \) and \( \TEDRm \).

Suppose a collection of centers \( ( \xv_{j} )_{j \in \JJJ} \), and 
the sets of directions \( (\Sph_{j})_{j \in \JJJ} \) have been fixed.
% 
Further, let 
\( \Imav_{j} = \bigl( \Ima_{j,\sph} \bigr)_{\sph \in \Sph_{j}} \),
\( \Sv_{j} = \bigl( S_{j,\sph} \bigr)_{\sph \in \Sph_{j}} \),
\( \UV_{j} = (\Uv_{\jm,\sph} \, , \sph \in \Sph_{j})^{\T} \),
be computed as in \eqref{ytd7e7ydwuhd5eythdm} with 
\( \weight_{ij} = \Kern(\| \TEDR_{0} (\Xv_{i} - \xv_{j}) \|^{2}) \)
for \( \TEDR_{0} = \hiso_{0}^{-1} \Id_{\dimd} \) and a starting bandwidth \( \hiso_{0} \).
%
Now, for any \( \xv_{\jl} \in (\xv_{j}) \), consider a localized version of
\eqref{ju7fudeju74334fhfd}:
\begin{EQA}[rcl]
%	\prmtv^{(0)} 
	&& \nquad
	\bigl\{ \PEDR_{\jl} \, , (\fcn_{\jjl},\fclv_{\jjl})_{j \in \JJJ} \bigr\}
	= 
	\arginf_{\PEDR \, , (\fcn_{j},\fclv_{j})_{j \in \JJJ}} 
	\sum_{j \in \JJJ} 
	\bigl\| \Imav_{j} - \fcn_{j} \Sv_{j} - \UV_{j} \PEDR^{\T} \fclv_{j} \bigr\|^{2}
	\weightM_{\jjl}
%	\Kern\bigl( \| \TEDRm(\xv_{j} - \xv_{\jl}) \|^{2} \bigr)
%	+ \lambda_{\PEDR} \pent(\BEDR)
	\, ,
	\qquad \quad
\label{kjhviir8786tytcdeefm}
\end{EQA}
where \( \weightM_{\jjl} = \Kern(\| \TEDRm_{0} (\xv_{j} - \xv_{\jl}) \|^{2}) \) for
\( \TEDRm_{0} = \hisom_{0}^{-1} \Id_{\dimd} \) and \( \hisom_{0} \gg \hiso_{0} \).
The only difference between this problem and \eqref{ju7fudeju74334fhfd} is in the multiplier (localizing weight)
\( \weightM_{\jjl} \)
%\( \Kern\bigl( \| \TEDRm(\xv_{j} - \xv_{\jl}) \|^{2} \bigr) \)
for each entry \( \bigl\| \Imav_{j} - \fcn_{j} \Sv_{j} - \UV_{j} \PEDR^{\T} \fclv_{j} \bigr\|^{2} \).
Such reweighting restricts the multi-index approximation to a vicinity of \( \xv_{\jl} \) given by 
such weights \( \weightM_{\jjl} \). 
%the tensor \( \TEDRm \).
%A solution can be found by alternating, exactly as in the multi-index case.
%Also, define
%\begin{EQA}
%	\JEDR_{\jl}
%	&=&
%	\BEDR_{\jl}^{\T} 
%		\biggl( \sum_{j \in \JJJ} \fclv_{\jjl} \, \fclv_{\jjl}^{\T} \, \weightM_{\jjl}
%%		\Kern\bigl( \| \TEDRm(\xv_{j} - \xv_{\jl}) \|^{2} \bigr)
%		\biggr) 
%	\BEDR_{\jl}
%	\, .
%\label{uhehdisnckncw9ud9}
%\end{EQA}
%%
%SVD-decomposition of \( \JEDR_{\jl} \) delivers the mapping \( \PEDR_{\jl} \) describing 
%the manifold structure at \( \xv_{\jl} \).
This local problem should be solved for any point from the family \( (\xv_{j}) \)
by alternating, exactly as in the multi-index case.
Note, however, that a minor change of \( \xv_{\jl} \) does not dramatically change the obtained solution.
Therefore, one can use a solution at one point as a starting guess for another close point.
This guarantees a fast convergence of the alternating procedure. 
All local problems \eqref{kjhviir8786tytcdeefm} can be merged into one.
This leads to minimization of
\begin{EQA}[c]
%	&& \nquad
%	(\BEDR_{\jl} \, , (\fcn_{\jjl},\fclv_{\jjl})_{j \in \JJJ} )_{\jl \in \JJJ}
%	\\
%	&=& 
%	\arginf_{(\BEDR_{j} \, , \fcn_{\jjl},\fclv_{\jjl})_{j \in \JJJ}} 
	\sum_{\jl \in \JJJ} \sum_{j \in \JJJ}
	\bigl\| \Imav_{j} - \fcn_{\jjl} \Sv_{j} - \UV_{j} \PEDR_{\jl}^{\T} \fclv_{\jjl} \bigr\|^{2} \weightM_{\jjl}
%	\Kern\bigl( \| \TEDRm(\xv_{j} - \xv_{\jl}) \|^{2} \bigr)
%	+ \lambda_{\PEDR} \pent(\BEDR)
	\, ,
	\qquad \quad
\label{kjhviir8786tytcdeefde}
\end{EQA}
over \( \{ \PEDR_{\jl} \, , (\fcn_{\jjl},\fclv_{\jjl})_{j \in \JJJ} \}_{\jl \in \JJJ} \).
%
%
The choice of the family \( (\PEDR_{\jl}) \) can be stabilized and synchronized by the penalty
\begin{EQA}[c]
	\lambda_{\PEDR} \, \pent_{\PEDR}(\PEDR)
	\eqdef
	\lambda_{\PEDR} \sum_{\jl \in \JJJ} \sum_{j \in \JJJ} \| \PEDR_{j} - \PEDR_{\jl} \|_{\Fr}^{2} 
	\; \weightM_{\jjl}
%	\Kern\bigl( \| \TEDRm_{j} (\xv_{j} - \xv_{\jl}) \|^{2} \bigr)
	\, .
\end{EQA}
%where the tensor \( \TEDRm_{j} \) describes the local manifold vicinity of the point \( \xv_{j} \).





\Subsection{Implementation} %. A flat manifold}
\label{SEDRmanproc}
%The multi-index procedure from Section~\ref{SEDRmiproc} can be extended to the case of manifold 
%hypothesis in a natural way.

Fix the tuning parameters (with defaut values in parenthesis)
\( \dimm (=1,2,3) \), 
\( \nmin (= 10) \), 
\( \nJJJ (= n/2) \), 
\( \nman (= \nJJJ/10) \),
\( \nSph (= \nmin \dimm) \),
\( a (= 2^{1/\dimm}) \), 
\( \lambda (= \nmin \dimm) \),
\( \lambda_{\PEDR} (=\nmin) \).

\paragraph{Step 0. Initialization.}
\begin{mylist}
\item
Fix the set \( (\xv_{j})_{j \in \JJJ} \).

\item
Fix \( \hiso_{0} \) and \( \hisom_{0} \) as the smallest values ensuring
\begin{EQA}[c]
	\frac{1}{\nJJJ} \sum_{j \in \JJJ} 
	\sumi \Kern\bigl( \hiso_{0}^{-2} \| \Xv_{i} - \xv_{j} \|^{2} \bigr)
	\geq 
	\nmin
	\, ,
	\quad
	\frac{1}{\nJJJ} \sum_{\jl \in \JJJ} 
	\sum_{j \in \JJJ} \Kern\bigl( \hisom_{0}^{-2} \| \xv_{j} - \xv_{\jl} \|^{2} \bigr)
	\geq 
	\nman
	\, .
\end{EQA}
%Also, set \( \haniso_{0} = \hiso_{0} \);

\item
For each \( j \in \JJJ \), generate the set \( \Sph_{j} \) of unit vectors \( \sph \) in \( \R^{\dimd} \)
with \( |\Sph_{j}| = \nSph \).

\item 
For each \( j \in \JJJ \) and \( \sph \in \Sph_{j} \), compute  
\( \Ima_{j,\sph}^{(0)} \, , \, S_{j,\sph}^{(0)} \, , \, \Uv_{\jm,\sph}^{(0)} \)
by \eqref{ytd7e7ydwuhd5eythdm} with 
\begin{EQA}[c]
	\weight_{ij} 
	= 
	\weight_{ij}^{(0)}
	= 
	\Kern\bigl( \hiso_{0}^{-2} \| \Xv_{i} - \xv_{j} \|^{2} \bigr) 
	\, .
\end{EQA}

\item 
With \( \Imav_{j}^{(0)} = \bigl( \Ima_{j,\sph}^{(0)} \bigr)_{\sph \in \Sph_{j}} \),
\( \Sv_{j}^{(0)} = \bigl( S_{j,\sph}^{(0)} \bigr)_{\sph \in \Sph_{j}} \),
\( \UV_{j}^{(0)} = (\Uv_{\jm,\sph}^{(0)} \, , \sph \in \Sph_{j})^{\T} \),
solve for every \( \jl \in \JJJ \) by alternating procedure from Section~\ref{SaltMI}
\begin{EQA}[rcl]
	&& \nquad
	\{ \PEDR_{\jl}^{(0)} \, , (\fcn_{\jjl}^{(0)},\fclv_{\jjl}^{(0)})_{j \in \JJJ} \} 
	\\
	&=& 
	\arginf_{\PEDR \, , (\fcn_{j},\fclv_{j})_{j \in \JJJ}} 
	\sum_{j \in \JJJ} 
	\bigl\| \Imav_{j}^{(0)} - \fcn_{j} \Sv_{j}^{(0)} - \UV_{j}^{(0)} \PEDR^{\T} \fclv_{j} \bigr\|^{2} \;
	\Kern\bigl( \hisom_{0}^{-2} \| \xv_{j} - \xv_{\jl} \|^{2} \bigr)
%	+ \lambda_{\PEDR} \pent(\BEDR)
	\, .
	\qquad \quad
\label{kjhviir8786tytcdeef}
\end{EQA}
The procedure also delivers \( \{ \JEDR_{\jl} \}_{\jl \in \JJJ} \).

\item 
Set \( k=1 \).
\end{mylist}


\paragraph{Iterations.}
For \( k \geq 1 \), let the last-step solution 
\( (\PEDR_{\jl}^{(k-1)} \, , \JEDR_{\jl}^{(k-1)} )_{\jl \in \JJJ} \) be available.
\begin{mylist}

\item
Update \( \hiso_{k} = \hiso_{k-1}/a \) and \( \hisom_{k} = \hisom_{k-1}/a \).

\item
Define \( \TEDR_{j}^{(k)} \) and \( \TEDRm_{\jl}^{(k)} \) by
\begin{EQA}[rcl]
	(\TEDR_{j}^{(k)})^{2}
	&=&
	\frac{\aniso_{k}^{2}}{\hiso_{k}^{2}} (\Id_{\dimd} - {\PEDR_{j}^{(k-1)}}^{\T} \PEDR_{j}^{(k-1)}) 
	+ \frac{1}{\hiso_{k}^{2}} \, \JEDR_{j}^{(k-1)}
	\, ,
	\qquad
	j \in \JJJ
	\, ,
	\\
	(\TEDRm_{\jl}^{(k)})^{2}
	&=&
	\frac{\anisom_{k}^{2}}{\hisom_{k}^{2}} (\Id_{\dimd} - {\PEDR_{\jl}^{(k-1)}}^{\T} \PEDR_{\jl}^{(k-1)}) 
	+ \frac{1}{\hisom_{k}^{2}} \, \JEDR_{\jl}^{(k-1)}
	\, ,
	\qquad
	\jl \in \JJJ
	\, ,
\end{EQA}
where the anisotropy indices \( \aniso_{k} \) and \( \anisom_{k} \) are the largest values ensuring
\begin{EQA}[rcl]
	\frac{1}{\nJJJ} \sum_{j \in \JJJ} \sumi \Kern\bigl( \| \TEDR_{j}^{(k)} (\Xv_{i} - \xv_{j}) \|^{2} \bigr)
	& \geq &
	\nmin
	\, ,
	\\
	\frac{1}{\nJJJ} \sum_{\jl \in \JJJ} \sum_{j \in \JJJ} \Kern\bigl( \| \TEDRm_{\jl}^{(k)} (\xv_{j} - \xv_{\jl}) \|^{2} \bigr)
	& \geq &
	\nman
	\, .
\label{jhgiverkytde4w34etycjamM}
\end{EQA}

\item
For each \( j \in \JJJ \), 
update the sets of directions \( \Sph_{j} \) by drawing \( \sph \eqd \gaussv_{\TEDR_{k}}/\| \gaussv_{\TEDR_{k}} \| \)
for \( \gaussv_{\TEDR_{k}} \sim \ND(0,\TEDR_{k}^{2}) \).

\item 
For each \( j \in \JJJ \) and \( \sph \in \Sph_{j} \), recompute
\( \Ima_{j,\sph}^{(k)} \), \( S_{j,\sph}^{(k)} \), and \( \Uv_{\jm,\sph}^{(k)} \) by \eqref{ytd7e7ydwuhd5eythdm} 
with %the weights  
\begin{EQA}[rcl]
	\weight_{ij} 
	= 
	\weight_{ij}^{(k)} 
	&=& 
	\Kern\bigl( \| \TEDR_{j}^{(k)} (\Xv_{i} - \xv_{j}) \|^{2} \bigr) 
	\, ,
	\qquad
	i = 1,\ldots,n
	\, ,
	\quad
	j \in \JJJ
	\, .
\end{EQA}
%With %\( \BEDR = (\BEDR_{j} \, , j \in \JJJ) \) for 
%\( \BEDR_{j} = (\betav_{j,1}, \ldots, \betav_{j,\dimm})^{\T} \in \R^{\dimm \times \dimd} \), 

\item
Solve by alternating
\begin{EQA}[rcl]
	&& \nquad
	\{ \PEDR_{\jl}^{(k)} \, , (\fcn_{\jjl}^{(k)},\fclv_{\jjl}^{(k)})_{j \in \JJJ} \}_{\jl \in \JJJ}
	\\
	&=& 
	\arginf_{\{ \BEDR_{\jl} \, , (\fcn_{\jjl},\fclv_{\jjl})_{j \in \JJJ} \}_{\jl \in \JJJ}} 
	\biggl\{ \sum_{\jl \in \JJJ} \sum_{j \in \JJJ} 
		\bigl\| \Imav_{j}^{(k)} - \fcn_{\jjl} \Sv_{j}^{(k)} - \UV_{j}^{(k)} \PEDR^{\T} \fclv_{\jjl} \bigr\|^{2}
		\Kern\bigl( \| \TEDRm_{\jl}^{(k)} (\xv_{j} - \xv_{\jl}) \|^{2} \bigr)
	\\
	&& \qquad
		+ \, \lambda \| \PEDR - \PEDR^{(k-1)} \|_{\Fr}^{2} + \lambda_{\PEDR} \pent(\PEDR)
	\biggr\}
	\, ,
\label{kjhviir8786tytcdeefk}
\end{EQA}
where
\begin{EQA}[c]
	\pent(\PEDR)
	\eqdef
	\sum_{\jl \in \JJJ} \sum_{j \in \JJJ} \| \PEDR_{j} - \PEDR_{\jl} \|_{\Fr}^{2} 
	\;
	\Kern\bigl( \| \TEDRm_{\jl}^{(k)} (\xv_{j} - \xv_{\jl}) \|^{2} \bigr)
	\, .
\end{EQA}

\item
Compute for every \( \jl \in \JJJ \)
\begin{EQA}[rcl]
	\JEDR_{\jl}^{(k)}
	& \eqdef &
	{\PEDR_{\jl}^{(k)}}^{\T} \biggl\{
		\sum_{j \in \JJJ} {\fclv}_{\jjl}^{(k)} \, ({\fclv}_{\jjl}^{(k)})^{\T} 
		\Kern\bigl( \| \TEDRm_{\jl} (\xv_{j} - \xv_{\jl}) \|^{2} \bigr)
	\biggr\} \PEDR_{\jl}^{(k)}
\end{EQA}
and its SVD
\begin{EQA}[c]
	\JEDR_{\jl}^{(k)}
	=
	{\PEDR_{\jl}^{(k)}}^{\T} \Lambda_{\jl}^{(k)} \, \PEDR_{\jl}^{(k)}
	\, ,
	\qquad
	\PEDR_{\jl}^{(k)} \, {\PEDR_{\jl}^{(k)}}^{\T}
	= 
	\Id_{\dimm}
	\, ,
	\quad
	\Lambda_{\jl}^{(k)} = \diag(\lambda_{1} \geq \ldots \geq \lambda_{\dimm})
	\, .
\end{EQA}

\item
If \( \hiso_{k}/a \geq \hiso_{\min} \), increase \( k \) by one and loop.

\item
Otherwise, terminate.
Output \( \PEDR_{\jl}^{(k)} \), \( \JEDR_{\jl}^{(k)} \), 
\( (\fcn_{\jjl}^{(k)},\fclv_{\jjl}^{(k)})_{j \in \JJJ} \}_{\jl \in \JJJ} \).
\end{mylist}



%%%%%%%%%%%%%%%%%%%%%%%%%%%%%%%  new section  %%%%%%%%%%%%%%%%%%%%%%%%%%%%%%%
\Section{Average derivative procedure}
\label{Saverder}

%One limitation of the proposed approach is its memory requirements: the \( (\dimd+1) \)-matrix \( \KSe_{j} \)
%has to be stored at every point \( \xv_{j} \).
%
This section provides a light modification of the method 
which is focused on estimation of the index subspace \( \EDR \).


\Subsection{Single-index}
\label{Saverdersi}
The single-index procedure from Section~\ref{Ssingind} recovers 
as the single-index vector \( \betav \), as the coefficients \( \fcn_{j}, \fcl_{j} \)
describing the value and slope coefficients of the function \( \fs \) at the point \( \xv_{j} \).
The average derivative approach focuses on estimation of directional derivatives. 
We use the notations of Section~\ref{Ssingind}.

Given a tensor \( \TEDR \), let \( \weight_{ij} = \Kern\bigl( \| \TEDR(\Xv_{i} - \xv_{j}) \|^{2} \bigr) \)
for all \( i=1,\ldots,n \) and \( j \in \JJJ \).
Also define the local mean
\begin{EQA}[c]
	\Xbar_{j}
	=
	\sumi \Xv_{i} \, \weight_{ij}
	\, ,
	\qquad
	j \in \JJJ
	\, .
\label{uyhvvy883hierwfewAVE}
\end{EQA}
%The main difference with the approach of Section~\ref{Ssingind} is the use of directional averaging. 
For each point \( \xv_{j} \), suppose a set \( \Sph_{j} \) of directions (unit vectors) \( \sph \in \R^{\dimd} \) 
to be given.
Compute for every \( \sph \in \Sph_{j} \)
\begin{EQ}[rcl]
	\Ima_{j,\sph}
	& \eqdef &
	\sumi Y_{i} \, \langle \Xv_{i} - \Xbar_{j} \, , \sph \rangle \, \weight_{ij}
	\, ,
	\\
	\Uv_{\jm,\sph}
	& \eqdef &
	\sumi (\Xv_{i} - \Xbar_{j}) \, \langle \Xv_{i} - \Xbar_{j} \, ,\sph \rangle  \, \weight_{ij}
	\in \R^{\dimd}
	\, .
	\qquad
\label{juu8hjfioee433erhAVE}
\end{EQ}
The model equation \eqref{YifXieiedrSI} yields
\begin{EQA}[c]
	\E \Ima_{j,\sph}
	=
	\sumi \fs(\betav^{\T} \Xv_{i}) \, \langle \Xv_{i} - \Xbar_{j} \, ,\sph \rangle \, \weight_{ij}
	\, .
\end{EQA}
Moreover, an approximation \( \fs(t) \approx \fcn_{j} + \fcl_{j} (t - t_{j}) \) implies by definition of \( \Xbar_{j} \)
\begin{EQA}[rcl]
	\E \Ima_{j,\sph}
	& \approx &
	\sumi \bigl\{ \fcn_{j} + \fcl_{j} \, \betav^{\T} (\Xv_{i} - \xv_{j}) \bigr\} \, \langle \Xv_{i} - \Xbar_{j} \, ,\sph \rangle \, \weight_{ij}
	\\
	&=&
	\sumi \fcl_{j} \, \betav^{\T} (\Xv_{i} - \Xbar_{j}) \, \langle \Xv_{i} - \Xbar_{j} \, ,\sph \rangle \, \weight_{ij}
%	=
%	\fcl_{j} \, \uv_{j}^{\T} \KSme_{j} \betav
	=
	\fcl_{j} \, \Uv_{\jm,\sph}^{\T} \, \betav
	\, .
\end{EQA}
%where
%\begin{EQA}[c]
%	\Uv_{\jm,\sph}
%	\eqdef
%	\sumi (\Xv_{i} - \Sv_{j}) \, \langle \Xv_{i} - \Sv_{j} \, ,\sph \rangle  \, \weight_{ij}
%	\, .
%	\qquad \quad
%\label{uyhvvy883hierwfewAVE}
%\end{EQA}
%Now we define
%\begin{EQA}
%	\LL_{\lambda}(\prmtv)
%	&=&
%	\frac{1}{2} \| \Imav - \imav \|^{2} 
%	+ \frac{1}{2} \sum_{j \in \JJJ} \Bigl\| \imav_{j} - \fcl_{j} \, \Uv_{\jm}^{\T} \betav \Bigr\|^{2} 
%	+ \frac{\lambda}{2} \| \betav - \betav_{0} \|^{2} \, .
%\label{tjjc7c6f5fAVE}
%\end{EQA}
With \( \Imav_{j} = (\Ima_{j,\sph})_{\sph \in \Sph_{j}} \in \R^{\nSph} \) and a \( \nSph \times \dimd \)-matrix 
\( \UV_{j} = (\Uv_{j,\sph})^{\T} \), this leads to the minimization of
\begin{EQA}[c]
	\sum_{j \in \JJJ} \sum_{\sph \in \Sph_{j}} \bigl( \Ima_{j,\sph} - \fcl_{j} \, \Uv_{\jm,\sph}^{\T} \, \betav \bigr)^{2} 
%	+ \lambda \| \betav - \betav_{0} \|^{2}
	=
	\sum_{j \in \JJJ} \bigl\| \Imav_{j} - \fcl_{j} \, \UV_{\jm} \, \betav \bigr\|^{2}
%	+ \lambda \| \betav - \betav_{0} \|^{2}
	\qquad
%	\to
%	\inf_{\betav,(\fcl_{j})_{j \in \JJJ}}
%	\, ,
\label{ikdfjoiefjoieuAVE}
\end{EQA}
over all unit vectors \( \betav \in \R^{\dimd} \) and \( (\fcl_{j})_{j \in \JJJ} \in \R^{\nSph} \), 
which can be solved by alternating minimization as for \eqref{ikdfjoiefjoieuvbjmr3e}.
For \( \betav^{(m-1)} \) fixed, each \( \fcl_{j} \) is estimated by
\begin{EQA}[c]
	\fcl_{j}^{(m)}
	=
	\frac{\langle \Imav_{j}, \UV_{\jm} \, \betav^{(m-1)} \rangle}{\| \UV_{\jm} \, \betav^{(m-1)} \|^{2}}
	\, ,
\end{EQA}
while for given \( (\fcl_{j}^{(m)})_{j \in \JJJ} \), the vector \( \betav^{(m)} \) can be estimated by
\begin{EQA}[c]
	\betav^{(m)}
	=
	\biggl( \sum_{j \in \JJJ} (\fcl_{j}^{(m)})^{2} \; \UV_{\jm}^{\T} \UV_{\jm} + \lambda \Id_{\dimp} \biggr)^{-1}
	\biggl( \sum_{j \in \JJJ} \fcl_{j}^{(m)} \, \UV_{\jm}^{\T} \Imav_{j} + \lambda \betav^{(m-1)} \biggr)
	\, .
\end{EQA}
{\color{blue} please check.}
Then normalize \( \betav^{(m)} \) to one and iterate until \( m > m_{\max} \).


\paragraph{Implementation. Step 0}
\begin{mylist}
\item
For each \( j \in \JJJ \),
fix a point \( \xv_{j} \in \R^{\dimd} \).

\item
Fix \( \hiso_{0} \) as the smallest value ensuring
\begin{EQA}[c]
	\frac{1}{\nJJJ} \sum_{j \in \JJJ} 
	\sumi \Kern\bigl( \hiso_{0}^{-2} \| \Xv_{i} - \xv_{j} \|^{2} \bigr)
%	=
%	\frac{1}{\nJJJ} \sum_{j \in \JJJ} \sumi \weight_{ij}
	\geq 
	\nmin
	\, ;
\end{EQA}

{\color{blue} please plot \( \hiso_{0} \)}

\item
Generate directions \( \sph \in \Sph_{j} \) uniformly on the unit sphere in \( \R^{\dimd} \).

\item 
for every \( j \in \JJJ \) and \( \sph \in \Sph_{j} \), compute \( \Ima_{j,\sph} \) and \( \Uv_{j,\sph} \) from 
\eqref{uyhvvy883hierwfewAVE} and \eqref{juu8hjfioee433erhAVE} 
with 
\begin{EQA}[c]
	\weight_{ij} = \Kern\bigl( \hiso_{0}^{-2} \| \Xv_{i} - \xv_{j} \|^{2} \bigr) 
	\, .
\end{EQA}
\item 
Find \( \betav_{0} \) by minimizing \eqref{ikdfjoiefjoieuAVE} using alternating minimization
with a random initialization.
%\item 
%%:
%rescale \( \betav_{0} \) by \( \| \betav_{0} \|^{-1} \) and \( (\fcl_{j}^{(0)}) \) by \( \| \betav_{0} \| \) to ensure 
%\( \| \betav_{0} \| = 1 \).
\end{mylist}

{\color{blue} check whether \( |\cos(\betav_{0},\betavs)| \) is significantly positive}

\paragraph{Implementation. Step \( k \)}
For \( k \geq 1 \), 
\begin{mylist}

\item 
update \( \hiso_{k} = \hiso_{k-1} / a \);
\item 
define \( \aniso_{k} \) by the condition
\begin{EQA}[c]
	\frac{1}{\nJJJ} \sum_{j \in \JJJ} \sumi 
		\Kern\Bigl( 
		\mfrac{\aniso_{k}^{2} \| \Xv_{i} - \xv_{j} \|^{2}}{\hiso_{k}^{2}} 
		+ \mfrac{\langle \Xv_{i} - \xv_{j}, \betav_{k-1} \rangle^{2}}{\hiso_{k}^{2}} 
	\Bigr)
	\geq 
	\nmin
	\, ,
\label{jhgiverkytde4w34etyAVE}
\end{EQA}
and fix 
\begin{EQA}[rcl]
	\TEDR_{k}^{2}
	&=&
	\hiso_{k}^{-2} \aniso_{k}^{2} \, \Id_{\dimd} + \hiso_{k}^{-2} \betav_{k-1} \betav_{k-1}^{\T}
	\, .
\label{iufhvu8kjbmnll475eAVE}
\end{EQA}

{\color{blue} plot \( \aniso_{k} \)}
\item
Renew the sets of directions \( \Sph_{j} \) using 
\( \sph \sim \gaussv_{\TEDR_{k}}/\| \gaussv_{\TEDR_{k}} \| \) for \( \gaussv_{\TEDR_{k}} \sim \ND(0,\TEDR_{k}^{2}) \).

\item 
For every \( j \in \JJJ \) and \( \sph \in \Sph_{j} \), compute \( \Ima_{j,\sph}^{(k)} \) and \( \Uv_{\jm,\sph}^{(k)} \) by \eqref{uyhvvy883hierwfewAVE} and \eqref{juu8hjfioee433erhAVE} 
with 
\begin{EQA}[c]
	\weight_{ij} 
	= 
	\weight_{ij}^{(k)}
	= 
	\Kern\bigl( \| \TEDR_{k} (\Xv_{i} - \xv_{j}) \|^{2} \bigr) 
	\, ;
\end{EQA}

\item 
Solve by alternating minimization with initialization \( \betav_{k-1} \):
\begin{EQA}[c]
	\{ \betav_{k},(\fcl_{j}^{(k)})_{j \in \JJJ} \} 
	=
	\arginf_{\betav,(\fcl_{j})_{j \in \JJJ}} 
	\sum_{j \in \JJJ} \bigl\| \Imav_{j}^{(k)} - \fcl_{j} \, \UV_{\jm}^{(k)} \betav \bigr\|^{2}
%	+ \lambda \| \betav - \betav_{k-1} \|^{2}
	\, .
\end{EQA}

{\color{blue} check whether \( |\cos(\betav_{0},\betavs)| \) grows to one with iterations}

%\item 
%%:
%Rescale \( \betav_{k} \) by \( \| \betav_{k} \|^{-1} \) and \( (\fcl_{j}^{(k)}) \) by \( \| \betav_{k} \| \) to ensure 
%\( \| \betav_{k} \| = 1 \).
\end{mylist}


\paragraph{Loop}
%
Increase \( k \) and loop while \( \hiso_{k} \geq \hiso_{\min}/a \).

\Subsection{Multiple index procedure}
\label{Saverdermi}
An extension to a multiple index procedure is straightforward.
Consider the e.d.r. regression problem
\begin{EQA}
	Y_{i}
	&=&
	\fs(\PEDR \Xv_{i}) + \eps_{i}
\label{YifXieiedrAVE}
\end{EQA}
where \( \Xv_{i} \in \R^{\dimd} \) for large \( \dimd \), 
\( \PEDR \) is an unknown element from the set \( \PROJ_{\dimd,\dimm} \) 
of orthogonal mappings from 
from \( \R^{\dimd} \) to \( \R^{\dimm} \) for \( \dimm \ll \dimd \), and
\( \fs \) is a function on \( \R^{\dimm} \).
The mapping \( \PEDR \) can be written as 
\begin{EQA}
	\PEDR \xv
	&=&
	(\betav_{1}^{\T} \xv) \betav_{1} + \ldots + (\betav_{\dimm}^{\T} \xv) \betav_{\dimm}
	\, ,
\label{ydgenb8jejAVE}
\end{EQA}
where \( \betav_{m}^{\T} \) are the orthonormal rows of \( \PEDR \).
Given a tensor \( \TEDR \), let the weights \( \weight_{ij} \) be defined as in \eqref{s7cuqkdwywdhbjwbfjwbdj}
for all \( i=1,\ldots,n \) and \( j \in \JJJ \).
The sums \( \Imav_{j} \) from \eqref{juu8hjfioee433erhAVE} can be expanded as
\begin{EQA}[rcl]
	\E \Ima_{j,\sph}
	& \approx &
	\sumi \biggl\{ \fcn_{j} + \sum_{\mm=1}^{\dimm} \fcl_{j,\mm} \, \betav_{\mm}^{\T} (\Xv_{i} - \xv_{j}) \biggr\} \, \langle \Xv_{i} - \Xbar_{j} \, ,\sph \rangle \, \weight_{ij}
	\\
	&=&
	\sum_{\mm=1}^{\dimm} \sumi \fcl_{j,\mm} \, \betav_{\mm}^{\T} (\Xv_{i} - \Xbar_{j}) \, \langle \Xv_{i} - \Xbar_{j} \, ,\sph \rangle \, \weight_{ij}
%	=
%	\fcl_{j} \, \uv_{j}^{\T} \KSme_{j} \betav
	=
	\Uv_{\jm,\sph}^{\T} \sum_{\mm=1}^{\dimm} \fcl_{j,\mm} \, \betav_{\mm}
	=
	\Uv_{\jm,\sph}^{\T} \PEDR^{\T} \fclv_{j}
	\, .
\end{EQA}
This suggests identifying the parameters \( (\fclv_{j})_{j \in \JJJ} \) and \( \PEDR \) by fitting the products 
\( \Uv_{\jm,\sph}^{\T} \PEDR^{\T} \fclv_{j} \) to \( \Ima_{j,\sph} \).
Therefore, each step of the procedure can be reduced to minimization of
\begin{EQA}[c]
	\sum_{j \in \JJJ} \sum_{\sph \in \Sph_{j}} 
	\bigl( \Ima_{j,\sph} - \Uv_{\jm,\sph}^{\T} \PEDR^{\T} \fclv_{j} \bigr)^{2}
%	+ \lambda \| \BEDR - \PEDR_{0} \|_{\Fr}^{2}
	=
	\sum_{j \in \JJJ}  
	\bigl\| \Imav_{j} - \UV_{\jm} \, \PEDR^{\T} \fclv_{j} \bigr\|^{2}
%	+ \lambda \| \BEDR - \PEDR_{0} \|_{\Fr}^{2}
	\, .
\label{ikdfjoiefvuheh3i}
\end{EQA}
We relax the constraint \( \PEDR \in \PROJ_{\dimd,\dimm} \) by 
replacing \( \PEDR \) with \( \BEDR \in \R^{\dimm \times \dimd} \), and
add stabilization penalty \( \lambda \| \BEDR - \PEDR_{0} \|_{\Fr}^{2} \)
for a preliminary guess \( \PEDR_{0} \) leading to the one-step problem
%
\begin{EQA}[c]
%	\sum_{j \in \JJJ} \sum_{\sph \in \Sph_{j}} 
%	\bigl( \Ima_{j,\sph} - \Uv_{\jm,\sph}^{\T} \BEDR^{\T} \fclv_{j} \bigr)^{2}
%	+ \lambda \| \BEDR - \PEDR_{0} \|_{\Fr}^{2}
%	=
	\inf_{\BEDR \, , \, (\fclv_{j})_{j \in \JJJ} }
	\sum_{j \in \JJJ}  
	\bigl\| \Imav_{j} - \UV_{\jm} \, \BEDR^{\T} \fclv_{j} \bigr\|^{2}
	+ \lambda \| \BEDR - \PEDR_{0} \|_{\Fr}^{2}
	\, 
\label{ikdfjoiefjoieuAVEmi}
\end{EQA}
over the set \( \BEDR \in \R^{\dimm \times \dimd} \), \( (\fclv_{j})_{j \in \JJJ} \in \R^{\dimm \nSph} \).
This problem can be effectively solved by alternating.

\begin{mylist}
\item
For the current solution \( \PEDR_{m-1} \) and each \( j \in \JJJ \), compute
\begin{EQA}[c]
	\fclv_{j}^{(m)}
	=
	\bigl( \PEDR_{m-1} \, \UV_{\jm}^{\T} \UV_{\jm} \, \PEDR_{m-1}^{\T} \bigr)^{-1} \PEDR_{m-1} \, \UV_{\jm}^{\T} \, \Imav_{j}
	\, .
\end{EQA}

\item
Then compute
\begin{EQA}[c]
	\BEDR_{m}
	=
	\arginf_{\BEDR} \biggl\{ \sum_{j \in \JJJ}  
	\bigl\| \Imav_{j} - \UV_{\jm} \, \BEDR^{\T} \fclv_{j}^{(m)} \bigr\|^{2}
	+ \lambda \| \BEDR - \PEDR_{m-1} \|_{\Fr}^{2} \biggr\}
	\, .
\end{EQA}
{\color{blue} please, provide a closed form solution}.
\item 
Compute
\begin{EQA}[rcl]
	\JEDR_{m}
	\eqdef
	\sum_{j \in \JJJ} (\BEDR_{m}^{\T} \, \fclv_{j}^{(m)}) \, \bigl( \BEDR_{m}^{\T} \, \fclv_{j}^{(m)} \bigr)^{\T}
	&=&
	\BEDR_{m}^{\T} \Bigl( \sum_{j \in \JJJ} {\fclv}_{j}^{(m)} \, ({\fclv}_{j}^{(m)})^{\T} \Bigr) \BEDR_{m}
	\, 
\end{EQA}
and its SVD
\begin{EQA}[c]
	\JEDR_{m}
	=
	\PEDR_{m}^{\T} \Lambda_{m} \, \PEDR_{m}
	\, ,
	\qquad
	\PEDR_{m} \, \PEDR_{m}^{\T}
	= 
	\Id_{\dimm}
	\, ,
	\quad
	\Lambda_{m} = \diag(\lambda_{1} \geq \ldots \geq \lambda_{\dimm})
	\, .
\end{EQA}

\item 
Iterate until \( m = m_{\max} \).
Output \( \PEDR_{m} \) and \( \JEDR_{m} \) for \( m = m_{\max} \).
\end{mylist}

The obtained solution can be used for structural adaptation as described in Section~\ref{SEDRsa}.
%
The whole procedure reads as follows. 


\paragraph{Implementation. Step 0}
\begin{mylist}
\item
Generate the set \( (\xv_{j} \, , j \in \JJJ) \).

\item
Fix \( \hiso_{0} \) as the smallest value ensuring
\begin{EQA}[c]
	\frac{1}{\nJJJ} \sum_{j \in \JJJ} 
	\sumi \Kern\bigl( \hiso_{0}^{-2} \| \Xv_{i} - \xv_{j} \|^{2} \bigr)
%	=
%	\frac{1}{\nJJJ} \sum_{j \in \JJJ} \sumi \weight_{ij}
	\geq 
	\nmin
	\, .
\end{EQA}

\item
For each \( j \in \JJJ \), generate the set \( \Sph_{j} \) of \( \nSph \) unit vectors \( \sph \) in \( \R^{\dimd} \).

\item 
For every \( j \in \JJJ \) and \( \sph \in \Sph_{j} \), compute \( \Ima_{j,\sph} \) and \( \Uv_{j,\sph} \) from 
\eqref{uyhvvy883hierwfewAVE} and \eqref{juu8hjfioee433erhAVE} 
with 
\begin{EQA}[c]
	\weight_{ij} 
	= 
	\weight_{ij}^{(0)}
	= 
	\Kern\bigl( \hiso_{0}^{-2} \| \Xv_{i} - \xv_{j} \|^{2} \bigr) 
	\, .
\end{EQA}
\item 
Find by alternating optimization with a random start \( \PEDR^{(0)} \)
\begin{EQA}[c]
	\{ \PEDR_{0} \, , (\fclv_{j}^{(0)})_{j \in \JJJ} \} 
	= 
	\arginf_{\PEDR \, , (\fclv_{j})_{j \in \JJJ}} 
	\sum_{j \in \JJJ} \bigl\| \Imav_{j} - \UV_{\jm} \, \PEDR^{\T} \, \fclv_{j} \bigr\|^{2}
	\, .
\end{EQA} 
%Here, the solution is not unique; any minimizer can be used.  
%from \eqref{gtgyuy6ew5dtde32322jtmi by alternating minimization.

\end{mylist}


\paragraph{Implementation. Step \( k \)}

\begin{mylist}
\item 
%:
Increase \( k \) by one. 

\item
Update \( \hiso_{k} = \hiso_{k-1} / a \).

\item
Define \( \TEDR_{k} \) by
\begin{EQA}[c]
	\TEDR_{k}^{2}
	\eqdef
	\hiso_{k}^{-2} \aniso_{k}^{2} \, \Id_{\dimd} + \hiso_{k}^{-2} \, \JEDR_{k-1}
	\quad
	\text{or }
	\quad
	\hiso_{k}^{-2} \aniso_{k}^{2} (\Id_{\dimd} - \PEDR_{k-1}^{\T} \PEDR_{k-1}) + \hiso_{k}^{-2} \, \JEDR_{k-1}
	\, ,
\end{EQA}
%{\color{blue} please, implement and compare both options}
%\\
where \( \aniso_{k} \) is the smallest value ensuring 
\begin{EQA}[c]
	\frac{1}{\nJJJ} \sum_{j \in \JJJ} \sumi 
		\Kern\bigl( \| \TEDR_{k} (\Xv_{i} - \xv_{j}) \|^{2} \bigr)
%	=
%	\frac{1}{\nJJJ} \sum_{j \in \JJJ} \sumi \weight_{ij}
	\geq 
	\nmin
	\, .
%\label{jhgiverkytde4w34etycjam}
\end{EQA}

\item
For each \( j \in \JJJ \), renew the sets of directions \( \Sph_{j} \) using 
\( \sph \sim \gaussv_{\TEDR_{k}}/\| \gaussv_{\TEDR_{k}} \| \) for \( \gaussv_{\TEDR_{k}} \sim \ND(0,\TEDR_{k}^{2}) \).

\item 
For every \( j \in \JJJ \) and \( \sph \in \Sph_{j} \), compute \( \Ima_{j,\sph}^{(k)} \) and \( \Uv_{\jm,\sph}^{(k)} \) by \eqref{uyhvvy883hierwfewAVE} and \eqref{juu8hjfioee433erhAVE} 
with 
\begin{EQA}[c]
	\weight_{ij} 
	= 
	\weight_{ij}^{(k)}
	= 
	\Kern\bigl( \| \TEDR_{k} (\Xv_{i} - \xv_{j}) \|^{2} \bigr) 
	\, .
\end{EQA}

\item 
Solve by alternating minimization starting from \( \PEDR_{k-1} \)
\begin{EQA}[c]
	\{ \PEDR_{k},(\fclv_{j}^{(k)})_{j \in \JJJ} \} 
	=
	\arginf_{\PEDR,(\fclv_{j})_{j \in \JJJ}} 
	\sum_{j \in \JJJ} \bigl\| \Imav_{j}^{(k)} - \UV_{\jm}^{(k)} \, \PEDR^{\T} \fclv_{j} \bigr\|^{2}
%	+ \lambda \| \BEDR - \BEDR_{k-1} \|^{2}
	\, .
\end{EQA}

%\item
%Compute the Jacobian matrix
%\begin{EQA}[rcl]
%	\JEDR_{k}
%	\eqdef
%	\sum_{j \in \JJJ} (\BEDR_{k}^{\T} {\fclv}_{j}^{(k)}) \, \bigl( \BEDR_{k}^{\T} {\fclv}_{j}^{(k)} \bigr)^{\T}
%	&=&
%	\BEDR_{k}^{\T} \Bigl( \sum_{j \in \JJJ} {\fclv}_{j}^{(k)} \, ({\fclv}_{j}^{(k)})^{\T} \Bigr) \BEDR_{k}
%	\, 
%\end{EQA}
%and its SVD
%\begin{EQA}[c]
%	\JEDR_{k}
%	=
%	\PEDR_{k}^{\T} \Lambda_{k} \, \PEDR_{k}
%	\, ,
%	\qquad
%	\PEDR_{k} \, \PEDR_{k}^{\T}
%	= 
%	\Id_{\dimm}
%	\, ,
%	\quad
%	\Lambda_{k} = \diag(\lambda_{1} \geq \ldots \geq \lambda_{\dimm})
%	\, .
%\end{EQA}

\item
Loop until \( \haniso_{k} < \haniso_{\min}/a \).

\end{mylist}


